\chapter{量子微分方程}\label{chap:quantum-differential-equations}\index{量子微分方程|hyperpage}
\SourcePage{324}
本章将介绍 Gromov--Witten 不变量的一种推广，称为\emph{引力关联数}，并用它们定义\emph{量子微分方程}。在\hyperref[sec:gravitational-correlators]{第~\ref*{sec:gravitational-correlators} 节}中，定义引力关联数并研究它们的性质。\hyperref[sec:giventals-connection]{第~\ref*{sec:giventals-connection} 节}随后利用引力关联数，显式描述上同调丛相对于\hyperref[chap:quantum-cohomology]{第八章}所构造的 A 模型联络的平坦截面。借助这些平坦截面，还将在\hyperref[sec:quantum-cohomology-relations]{第~\ref*{sec:quantum-cohomology-relations} 节}中定义一个取值于 $H^*(X)$、定义在 $H^*(X)$ 上的形式函数 $J$。这个函数可看作由某些较为重要的引力关联数组成的生成函数。特别地，Calabi--Yau 三维簇的 A 模型关联函数可以完全从 $J$ 中提取出来。然后，将量子微分算子定义为零化 $J$ 的微分算子。将证明，这些算子给出量子上同调中的关系，这也说明了 $J$ 本身的重要性。此函数在 Givental 对镜像定理的研究方法中同样起着核心作用；这一方法将在\hyperref[sec:givental-mirror-approach]{第~\ref*{sec:givental-mirror-approach} 节}讨论。

\section{引力关联数}\label{sec:gravitational-correlators}\index{引力关联数|hyperpage}
本节讨论引力关联数及其性质。它们是光滑代数簇的不变量，并包含\hyperref[chap:gromov-witten]{第七章}所研究的 Gromov--Witten 不变量。在\hyperref[subsec:gravitational-correlator-definition]{第~\ref*{subsec:gravitational-correlator-definition} 小节}中定义引力关联数，在\hyperref[subsec:gravitational-correlator-properties]{第~\ref*{subsec:gravitational-correlator-properties} 小节}中研究它们的性质；最后给出一系列例子。

\subsection{关联数的定义}\label{subsec:gravitational-correlator-definition}
定义引力关联数之前，需要回顾稳定映射的一些事实。给定光滑代数簇 $X$ 和 $\beta\in H_2(X,\mathbb Z)$，由\hyperref[subsec:stable-maps]{第~\ref*{subsec:stable-maps} 小节}有模叠 $\overline{\mathcal M}_{g,n}(X,\beta)$，及其对应的粗模空间 $\overline M_{g,n}(X,\beta)$。还有赋值映射
\[
e_i:\overline{\mathcal M}_{g,n}(X,\beta)\longrightarrow X,
\qquad i=1,\ldots,n,
\]
由 $e_i(f,p_1,\ldots,p_n)=f(p_i)$ 给出，以及映射
\[
\pi_{n+1}:\overline{\mathcal M}_{g,n+1}(X,\beta)
\longrightarrow\overline{\mathcal M}_{g,n}(X,\beta),
\]
它通过忘却最后一个标记点并收缩不稳定分支来定义。

由于 $\overline M_{g,n}(X,\beta)$ 不是 $\overline{\mathcal M}_{g,n}(X,\beta)$ 的细模空间，它不能作为一个万有曲线的参数空间。不过，我们断言，$\pi_{n+1}$ 是万有曲线在代数叠意义下的版本，而 $e_{n+1}:\overline{\mathcal M}_{g,n+1}(X,\beta)\to X$ 是万有稳定映射。为说明这一点，仍用 $e_{n+1}$、$\pi_{n+1}$ 表示粗模空间 $\overline M_{g,n+1}(X,\beta)$
\SourcePage{325}
和 $\overline M_{g,n}(X,\beta)$ 上对应的映射。对任何 $(f,C,p_1,\ldots,p_n)\in\overline M_{g,n}(X,\beta)$，有自然同构
\begin{equation}\label{eq:universal-stable-curve-coarse-fiber}
C/\operatorname{Aut}(f)\simeq\pi_{n+1}^{-1}(f),
\end{equation}
其中 $\operatorname{Aut}(f)$ 是稳定映射 $f$ 的有限自同构群。这个同构如下定义：若 $p\in C$ 是不同于 $p_1,\ldots,p_n$ 的光滑点，就将它映为 $(f:C\to X,p_1,\ldots,p_n,p)$。若 $p=p_i$，则将它映为稳定映射 $(f:C\cup\mathbb P^1\to X,p_1,\ldots,p_i',\ldots,p_n,p_{n+1}')$，其中 $\mathbb P^1$ 在 $p_i$ 处附着于 $C$，$p_i'$ 和 $p_{n+1}'$ 是 $\mathbb P^1$ 上两个互异且不同于附着点的点。最后，$p$ 为结点的情形留给读者处理。由于可将 $f$ 与 $e_{n+1}$ 在 $\pi_{n+1}^{-1}(f)$ 上的限制认同，我们的断言得到证明。在对应的代数叠表述中，自同构已被叠结构计入，所以在代数叠意义下，$\pi_{n+1}$ 就是万有曲线。

式~\eqref{eq:universal-stable-curve-coarse-fiber} 对 $\pi_{n+1}$ 纤维的描述表明，它具有典范截面 $s_i:\overline M_{g,n}(X,\beta)\to\overline M_{g,n+1}(X,\beta)$，定义为
\[
s_i(f,C,p_1,\ldots,p_n)
=(f:C\cup\mathbb P^1\to X,p_1,\ldots,p_i',\ldots,p_n,p_{n+1}').
\]
按式~\eqref{eq:universal-stable-curve-coarse-fiber} 的描述，此映射就是将 $(f,C,p_1,\ldots,p_n)$ 送到 $p_i\in C$ 的映射。

现在令 $\mathcal L_i$ 表示“第 $i$ 个标记点处的余切线”，即在稳定映射 $(f:C\to X,p_1,\ldots,p_n)$ 处以余切空间 $T_{p_i}^*C$ 为纤维的线丛。严格地说，令 $\omega_{n+1}$ 为 $\pi_{n+1}$ 的相对对偶化层，则 $\mathcal L_i=s_i^*\omega_{n+1}$，这里使用代数叠上的万有曲线及其截面。注意，$\mathcal L_i$ 是叠 $\overline{\mathcal M}_{g,n}(X,\beta)$ 上的线丛，却不一定是粗模空间 $\overline M_{g,n}(X,\beta)$ 上的局部自由层。除了线丛 $\mathcal L_i$ 和赋值映射 $e_i$，还需要\IndexLocation{indexloc:1406:325}虚拟基本类 $[\overline M_{g,n}(X,\beta)]^{\mathrm{virt}}$。由\hyperref[subsec:virtual-fundamental-class]{第~\ref*{subsec:virtual-fundamental-class} 小节}，其同调次数为 $2(1-g)(d-3)-2\int_\beta\omega_X+2n$，其中 $d$ 是 $X$ 的复维数，$\omega_X$ 为典范类。现在可以作出定义。
\begin{definition}\label{def:gravitational-correlator}
给定类 $\gamma_1,\ldots,\gamma_n\in H^*(X)$，以及非负整数 $d_i$（$i=1,\ldots,n$），不变量
\[
\langle\tau_{d_1}\gamma_1,\ldots,\tau_{d_n}\gamma_n\rangle_{g,\beta}
=\int_{[\overline M_{g,n}(X,\beta)]^{\mathrm{virt}}}
\prod_{i=1}^n\bigl(c_1(\mathcal L_i)^{d_i}\cup e_i^*(\gamma_i)\bigr)
\]
称为\emph{引力关联数}。
\end{definition}

注意，只有在 $\overline M_{g,n}(X,\beta)$ 有定义时，引力关联数才有定义。这会限制三元组 $(g,n,\beta)$ 的取值。

定义~\ref{def:gravitational-correlator} 中的积分将 $\mathcal L_i$ 当作模空间 $\overline M_{g,n}(X,\beta)$ 上的线丛。严格来说，这个积分应在叠 $\overline{\mathcal M}_{g,n}(X,\beta)$ 上求值。不过，如果认为被积式已推前到 $\overline M_{g,n}(X,\beta)$ 上，那么在粗模空间上求积分就完全有意义。这里以及本书余下部分，都采用这一约定。它在思想上类似于例题~\ref{ex:quintic-virtual-class} 中对丛 $\mathcal V_d$ 及其 Euler 类所采用的约定。

引力关联数这一名称来自物理学中的类似术语。如\hyperref[sec:nonlinear-sigma-models]{附录 B.2} 所述，非线性 sigma 模型将算子 $\mathcal O_\gamma$（作为 BRST 上同调类）赋予每个上同调类 $\gamma$。算子 $\mathcal O_\gamma$ 又产生其他物理算子 $\mathcal O_{i,\gamma}$，其中整数 $i\ge0$；这些算子称为\emph{引力后裔}。引力关联数与引力后裔
\SourcePage{326}
之间的关系由方程
\begin{equation}\label{eq:gravitational-physical-normalization}
\langle\tau_{d_1}\gamma_1,\ldots,\tau_{d_n}\gamma_n\rangle_{g,\beta}
=\frac{\langle\mathcal O_{d_1,\gamma_1},\ldots,\mathcal O_{d_n,\gamma_n}\rangle_{g,\beta}}
{\prod_{j=1}^n d_j!}
\end{equation}
给出。式~\eqref{eq:gravitational-physical-normalization} 的左端是定义~\ref{def:gravitational-correlator} 中的数学引力关联数。右端的表达式 $\langle\mathcal O_{d_1,\gamma_1},\ldots,\mathcal O_{d_n,\gamma_n}\rangle_{g,\beta}$ 是物理关联函数 $\langle\mathcal O_{d_1,\gamma_1},\ldots,\mathcal O_{d_n,\gamma_n}\rangle$ 的一部分：这一部分使用亏格为 $g$、满足 $f_*[\Sigma]=\beta$ 的映射 $f:\Sigma\to X$，以及引力后裔场来计算。\hyperref[sec:general-field-theories]{附录 B.1} 概述了物理关联函数的思想。式~\eqref{eq:gravitational-physical-normalization} 右端分母中的因子纯属约定。

一个重要观察是，Gromov--Witten 不变量恰好是所有 $d_i=0$ 的引力关联数。若将 $\tau_0\gamma$ 简写为 $\gamma$，则按\hyperref[chap:gromov-witten]{第七章}的记号，
\[
\langle\gamma_1,\ldots,\gamma_n\rangle_{g,\beta}
=\langle I_{g,n,\beta}\rangle(\gamma_1,\ldots,\gamma_n).
\]
本章对 Gromov--Witten 不变量使用较简便的记号 $\langle\gamma_1,\ldots,\gamma_n\rangle_{g,\beta}$。还经常将 $\tau_d1$ 简写为 $\tau_d$，其中 $1\in H^0(X)$ 是 $X$ 的基本类所对应的上同调单位。

引力关联数还有一些有用的组合，可产生形式函数。先回顾\hyperref[chap:quantum-cohomology]{第八章}的记号：取 $H^*(X,\mathbb Q)$ 的基 $T_0=1,T_1,\ldots,T_m$。又如\hyperref[subsec:gw-potential]{第~\ref*{subsec:gw-potential} 小节}中定义 Gromov--Witten 势时一样，令 $\gamma=\sum_i t_iT_i$。若 $\omega$ 是 $X$ 上的复化 Kähler 类，则定义\emph{亏格 $g$ 耦合}
\begin{equation}\label{eq:genus-g-gravitational-coupling}
\begin{aligned}
\langle\!\langle\tau_{d_1}\gamma_1,\ldots,\tau_{d_n}\gamma_n\rangle\!\rangle_g
=\sum_{k=0}^{\infty}\sum_\beta\frac1{k!}
\langle\tau_{d_1}\gamma_1,\ldots,\tau_{d_n}\gamma_n,
\underbrace{\gamma,\ldots,\gamma}_{k\text{ 次}}\rangle_{g,\beta}\,q^\beta,
\end{aligned}
\end{equation}
其中，如往常一样，$q^\beta=e^{2\pi i\int_\beta\omega}$。在不会混淆时，有时将右端求和中的引力关联数简写为
\[
\langle\tau_{d_1}\gamma_1,\ldots,\tau_{d_n}\gamma_n,(\gamma)^k\rangle_{g,\beta}.
\]

类似于\hyperref[subsec:formal-small-quantum]{第~\ref*{subsec:formal-small-quantum} 小节}和\hyperref[subsec:gw-potential]{第~\ref*{subsec:gw-potential} 小节}中的处理，可以把式~\eqref{eq:genus-g-gravitational-coupling} 的系数解释为属于某个环 $\widehat C$；视具体情形，可取
\[
\begin{gathered}
\mathbb C[[t_0,\ldots,t_m]],\quad R[[t_0,\ldots,t_m]],\\
\mathcal R[[t_0,\ldots,t_m]],\quad
\Lambda(\omega,\mathbb Q)[[t_0,\ldots,t_m]].
\end{gathered}
\]
因此，引力关联数及亏格 $g$ 耦合都是定义良好的形式对象。

用\emph{亏格 $g$ 引力 Gromov--Witten 势}重新表述这些内容，会更清楚。对 $d\ge0$ 和 $0\le j\le m$，引入超交换变量 $t_d^j$，使 $\deg t_d^j=\deg T_j$，其中 $t_0^j$ 就是通常的变量 $t_j$。随后修改 Gromov--Witten 势 $\Phi$ 的定义（定义~\ref{def:gw-potential}），得到亏格 $g$ 引力 Gromov--Witten 势 $\Phi_g^{\mathrm{grav}}$。具体地，将表达式 $\gamma=\sum_i t_iT_i$ 替换为形式表达式
\[
\gamma=\sum_{d=0}^{\infty}\sum_{j=0}^m t_d^j\tau_dT_j,
\]
然后用亏格 $g$ 映射代替亏格零映射，再如\hyperref[subsec:gw-potential]{第~\ref*{subsec:gw-potential} 小节}中那样展开。更准确地说，有如下定义。
\SourcePage{327}
\begin{definition}\label{def:gravitational-gw-potential}\index{Gromov--Witten势@Gromov--Witten 势!引力势|hyperpage}
设 $\omega$ 是光滑射影簇 $X$ 上的复化 Kähler 类。则\emph{亏格 $g$ 引力 Gromov--Witten 势}为形式和
\[
\Phi_g^{\mathrm{grav}}(\gamma)
=\sum_{n=0}^{\infty}\sum_{\beta\in H_2(X,\mathbb Z)}
\frac1{n!}\langle\gamma^n\rangle_{g,\beta}\,q^\beta,
\]
其中 $q^\beta=e^{2\pi i\int_\beta\omega}$。
\end{definition}
在上述求和中，只对使 $\overline M_{g,n}(X,\beta)$ 存在的 $(n,\beta)$ 求和。注意，
\begin{equation}\label{eq:gravitational-potential-derivatives}
\langle\!\langle\tau_{d_1}T_{i_1},\ldots,\tau_{d_n}T_{i_n}\rangle\!\rangle_g
=\left.\frac{\partial^n\Phi_g^{\mathrm{grav}}}
{\partial t_{d_1}^{i_1}\cdots\partial t_{d_n}^{i_n}}
\right|_{t_d^j=0\text{，所有 }d>0},
\end{equation}
这里，对于奇变量的偏导数，采用\hyperref[subsec:cohomology-supermanifold]{第~\ref*{subsec:cohomology-supermanifold} 小节}中的符号约定。特殊地，当 $g=0$ 且所有 $d_i=0$ 时，可写成
\[
\langle\!\langle T_{j_1},\ldots,T_{j_n}\rangle\!\rangle_0
=\frac{\partial^n\Phi}{\partial t_{j_1}\cdots\partial t_{j_n}},
\]
其中 $\Phi$ 是通常的 Gromov--Witten 势。

立即可知，大量子乘积（定义~\ref{def:big-quantum-product}）可改写为
\begin{equation}\label{eq:big-quantum-product-couplings}
T_i*T_j=\sum_k\langle\!\langle T_i,T_j,T_k\rangle\!\rangle_0\,T^k,
\end{equation}
这里利用 $g^{kl}$ 升指标，令 $T^k=\sum_l g^{kl}T_l$。如\hyperref[chap:quantum-cohomology]{第八章}所述，这个乘积通过 Dubrovin 形式体系由 Gromov--Witten 势 $\Phi$ 得到。若将 $\Phi$ 替换为 $\Phi_0^{\mathrm{grav}}$，便在 $H^*(X)$ 上得到一个新的结合乘积，定义为
\[
T_i*_gT_j=\sum_k\frac{\partial^3\Phi_0^{\mathrm{grav}}}
{\partial t_i\partial t_j\partial t_k}\,T^k.
\]
这就是\emph{引力量子乘积}。其中的系数可视具体情形属于以下某个环：
\[
\mathbb C[[t_d^j]],\qquad R[[t_d^j]],\qquad
\mathcal R[[t_d^j]],\qquad\Lambda(\omega,\mathbb Q)[[t_d^j]].
\]

\subsection{关联数的性质}\label{subsec:gravitational-correlator-properties}
刚刚定义的引力关联数具有与\hyperref[sec:gw-properties]{第~\ref*{sec:gw-properties} 节}中 Gromov--Witten 不变量的性质相似的性质。这里只陈述将要用到的那些。它们来自虚拟基本类的性质 \SourceCite{BM,Behrend}。
\begin{axiom*}{线性}
$\langle\tau_{d_1}\gamma_1,\ldots,\tau_{d_n}\gamma_n\rangle_{g,\beta}$ 对每个变元都是线性的。
\end{axiom*}

\begin{axiom*}{有效性}
对光滑射影簇 $X$，若 $\beta$ 不是有效类，则 $\langle\tau_{d_1}\gamma_1,\ldots,\tau_{d_n}\gamma_n\rangle_{g,\beta}=0$。
\end{axiom*}

\begin{axiom*}{次数}
对齐次类 $\gamma_i$，引力关联数
\[
\langle\tau_{d_1}\gamma_1,\ldots,\tau_{d_n}\gamma_n\rangle_{g,\beta}
\]
只有在
\[
\sum_{i=1}^n(2d_i+\deg\gamma_i)
=2(1-g)\dim X-2\int_\beta\omega_X+2(3g-3+n)
\]
时才可能非零。
\end{axiom*}
这容易由定义~\ref{def:gravitational-correlator} 和虚拟基本类的维数公式得到。

\SourcePage{328}
\begin{axiom*}{基本类}
如\hyperref[chap:gromov-witten]{第七章}中一样，$\pi_n:\overline M_{g,n}(X,\beta)\to\overline M_{g,n-1}(X,\beta)$ 是忘却最后一点并收缩不稳定分支的映射。若 $n+2g\ge4$，或 $\beta\ne0$ 且 $n\ge1$，则此映射有定义。令 $\mathcal L_i'$ 为 $\overline M_{g,n-1}(X,\beta)$ 上第 $i$ 个标记点处的余切线丛，则有恒等式
\begin{equation}\label{eq:cotangent-line-forgetting-boundary}
c_1(\mathcal L_i)=\pi_n^*c_1(\mathcal L_i')
+\widetilde D_{(i,n\mid1,\ldots,i-1,i+1,\ldots,n-1)},
\end{equation}
其中 $\widetilde D_{(i,n\mid1,\ldots,i-1,i+1,\ldots,n-1)}$ 表示如下稳定映射轨迹的闭包：源曲线有两个分支 $C_1,C_2$，$C_1$ 上有标记点 $\{i,n\}$，$C_2$ 上有标记点 $\{1,\ldots,i-1,i+1,\ldots,n-1\}$，而且 $f$ 在 $C_1$ 上为常值。又如式~\eqref{eq:virtual-class-forgetting-compatibility} 所述，$\overline M_{g,n}(X,\beta)$ 的虚拟基本类是 $\overline M_{g,n-1}(X,\beta)$ 的虚拟基本类的拉回。将这一事实、式~\eqref{eq:cotangent-line-forgetting-boundary} 和 \SourceCite{Witten3} 中的计算结合，得到
\[
\begin{aligned}
&\langle\tau_{d_1}\gamma_1,\ldots,\tau_{d_{n-1}}\gamma_{n-1},1\rangle_{g,\beta}\\
&\quad=\sum_{i=1}^{n-1}\langle\tau_{d_1}\gamma_1,\ldots,
\tau_{d_{i-1}}\gamma_{i-1},\tau_{d_i-1}\gamma_i,
\tau_{d_{i+1}}\gamma_{i+1},\ldots,\tau_{d_{n-1}}\gamma_{n-1}\rangle_{g,\beta},
\end{aligned}
\]
其中任何含有 $\tau_{-1}$ 的项均取为零。
\end{axiom*}

\begin{axiom*}{除子}
设 $D$ 为除子。式~\eqref{eq:cotangent-line-forgetting-boundary} 及虚拟基本类的相容性给出
\[
\begin{aligned}
&\langle\tau_{d_1}\gamma_1,\ldots,\tau_{d_{n-1}}\gamma_{n-1},D\rangle_{g,\beta}\\
&=\left(\int_\beta D\right)
\langle\tau_{d_1}\gamma_1,\ldots,\tau_{d_{n-1}}\gamma_{n-1}\rangle_{g,\beta}\\
&\quad+\sum_{j=1}^{n-1}\langle\tau_{d_1}\gamma_1,\ldots,
\tau_{d_{j-1}}\gamma_{j-1},\tau_{d_j-1}(D\cup\gamma_j),
\tau_{d_{j+1}}\gamma_{j+1},\ldots,\tau_{d_{n-1}}\gamma_{n-1}\rangle_{g,\beta}.
\end{aligned}
\]
如同基本类公理，这条公理在 $n+2g\ge4$，或 $\beta\ne0$ 且 $n\ge1$ 时成立。
\end{axiom*}

\begin{axiom*}{分裂}
为便于陈述这条公理，稍微改变观点，将 Gromov--Witten 类的定义推广，使它也包含引力项。沿用定义~\ref{def:gw-algebraic-invariants} 的记号，若 $2g+n\ge3$，$\gamma_1,\ldots,\gamma_n\in H^*(X,\mathbb Q)$，则称由
\begin{equation}\label{eq:gravitational-class-definition}
\begin{aligned}
I_{g,n,\beta}(\tau_{d_1}\gamma_1,\ldots,\tau_{d_n}\gamma_n)
=PD^{-1}p_{2*}\pi_*\left((p_1\circ\pi)^*(\gamma_1\otimes\cdots\otimes\gamma_n)
\cap\left(\left(\prod_{i=1}^n c_1(\mathcal L_i)^{d_i}\right)\cup\xi\right)\right)
\end{aligned}
\end{equation}
给出的 $H^*(\overline M_{g,n},\mathbb Q)$ 中的元素为\emph{引力类}。其中，$\xi$ 是虚拟基本类，
\[
PD:H^*(\overline M_{g,n},\mathbb Q)
\simeq H_{6g-6+2n-*}(\overline M_{g,n},\mathbb Q)
\]
是 Poincaré 对偶同构。这里仍采用先前的约定：先将括号内的表达式看作叠 $\overline{\mathcal M}_{g,n}(X,\beta)$ 上的类，再推前到粗模空间 $\overline M_{g,n}$。在这一记号下，注意
\begin{equation}\label{eq:gravitational-class-integral}
\langle\tau_{d_1}\gamma_1,\ldots,\tau_{d_n}\gamma_n\rangle_{g,\beta}
=\int_{\overline M_{g,n}}I_{g,n,\beta}(\tau_{d_1}\gamma_1,\ldots,\tau_{d_n}\gamma_n).
\end{equation}
如\hyperref[sec:gw-properties]{第~\ref*{sec:gw-properties} 节}中一样，考虑将两条带标记曲线的粘合点接在一起的映射
\[
\varphi:\overline M_{g_1,n_1+1}\times\overline M_{g_2,n_2+1}
\longrightarrow\overline M_{g,n},
\qquad g=g_1+g_2,\quad n=n_1+n_2.
\]
\SourcePage{329}
在稳定映射的源曲线上施行这种粘合，便形成新的稳定映射。分裂公理断言
\[
\begin{aligned}
&\varphi^*I_{g,n,\beta}(\tau_{d_1}\gamma_1,\ldots,\tau_{d_n}\gamma_n)\\
&=\sum_{i,\,\beta=\beta_1+\beta_2}
I_{g_1,n_1+1,\beta_1}(\tau_{d_1}\gamma_1,\ldots,\tau_{d_{n_1}}\gamma_{n_1},T_i)\otimes\\
&\hspace{5em}I_{g_2,n_2+1,\beta_2}(T^i,\tau_{d_{n_1+1}}\gamma_{n_1+1},\ldots,\tau_{d_n}\gamma_n).
\end{aligned}
\]
\end{axiom*}

\begin{axiom*}{伸缩}\index{伸缩公理|hyperpage}
最后一条公理在 Gromov--Witten 理论中没有对应物，它是恒等式
\[
\langle\tau_1,\tau_{d_1}\gamma_1,\ldots,\tau_{d_n}\gamma_n\rangle_{g,\beta}
=(2g-2+n)\langle\tau_{d_1}\gamma_1,\ldots,\tau_{d_n}\gamma_n\rangle_{g,\beta}.
\]
为解释这一点，将万有曲线 $C_{g,n}(X,\beta)\to\overline M_{g,n}(X,\beta)$ 与基本类公理中使用的映射 $\pi_{n+1}:\overline M_{g,n+1}(X,\beta)\to\overline M_{g,n}(X,\beta)$ 认同。按这一认同，新增标记点对应的余切线丛在 $\pi_{n+1}$ 的纤维上，即在稳定映射 $(f,C,p_1,\ldots,p_n)$ 的源曲线 $C$ 上，限制为允许在各 $p_i$ 处有至多一阶极点的对偶化微分层 $\omega_C(p_1+\cdots+p_n)$。于是，由虚拟基本类相对于 $\pi_{n+1}$ 的相容性，以及该层在每个纤维上的次数为 $2g-2+n$，得到伸缩公理。
\end{axiom*}

虽然引力关联数看上去比 Gromov--Witten 不变量复杂，但 \SourceCite{Dubrovin2} 已注意到，在亏格零情形中，可以用 Gromov--Witten 不变量 $\langle\gamma_1,\ldots,\gamma_n\rangle_{0,\beta}$ 计算所有亏格零引力关联数。Dubrovin 的证明使用除子公理和下面的引理~\ref{lem:genus-zero-topological-recursion}。这里简要介绍 R. Pandharipande 的另一种方法。

其思想是先增加三个辅助标记点 $p_{n+1}$、$p_{n+2}$ 和 $p_{n+3}$。接着，将引力关联数定义中的被积式拉回到模空间 $\overline M_{0,n+3}(X,\beta)$，再与下列类作杯积：
\[
\phi=e_{n+1}^*(H)\cup e_{n+2}^*(H)
\]
其中 $H$ 是 $X$ 上的丰沛类。忘却前 $n-1$ 个标记点和映射，便得到 $s:\overline M_{0,n+3}(X,\beta)\to\overline M_{0,4}$。然后将 $\phi$ 分别与式~\eqref{eq:wdvv-boundary-equivalence} 中线性等价关系
\[
D(n,n+1\mid n+2,n+3)\sim D(n,n+3\mid n+1,n+2)
\]
两端在 $s$ 下的拉回作杯积。使用除子公理、基本类公理和分裂公理，就得到用亏格零 Gromov--Witten 不变量递归计算亏格零引力关联数的方法。下面的例题~\ref{ex:projective-line-gravitational} 将概述这一方法。将引力关联数约化的另一种方法见 \SourceCite{KoM3}。

\subsection{关联数的例题}\label{subsec:gravitational-correlator-examples}
现在给出几个例子，说明引力关联数的计算。为方便起见，当 $X=\mathbb P^r$、$\beta=d[\ell]$，其中 $\ell\subset\mathbb P^r$ 是直线时，有时用 $\overline M_{g,n}(\mathbb P^r,d)$ 代替 $\overline M_{g,n}(\mathbb P^r,\beta)$，用 $\langle\tau_{d_1}\gamma_1,\ldots,\tau_{d_n}\gamma_n\rangle_{g,d}$ 代替 $\langle\tau_{d_1}\gamma_1,\ldots,\tau_{d_n}\gamma_n\rangle_{g,\beta}$。
\SourceExampleNumber{10.1.3}{1}
\begin{example}\label{ex:projective-line-gravitational}
将计算 $X=\mathbb P^1$ 的一些引力关联数。如\hyperref[subsec:moduli-to-gw]{第~\ref*{subsec:moduli-to-gw} 小节}所述，$\overline M_{0,n}(\mathbb P^1,d)$ 的维数为\IndexLocation{indexloc:1279:329}预期维数 $2d+n-2$，因此计算中不必引入虚拟基本类。

先计算次数为 $1$ 的不变量。注意，$\overline M_{0,1}(\mathbb P^1,1)\simeq\mathbb P^1$，同构将 $f:(C,p)\to\mathbb P^1$ 映为 $f(p)$。在此同构下，$\mathcal L_1$ 对应于 $\mathbb P^1$ 的余切丛，因此得到一点引力
\SourcePage{330}
关联数 $\langle\tau_1\rangle_{0,1}=-2$。令 $H$ 为 $\mathbb P^1$ 的超平面类，也就是一个点的类。则有 Gromov--Witten 不变量
\[
\langle H\rangle_{0,1}=\langle H,H,H\rangle_{0,1}
=\langle I_{0,3,1}\rangle(H,H,H)
=\langle I_{0,3,1}\rangle(pt,pt,pt)=1.
\]
第一个等号来自除子公理，而最后一个等号是显然的，因为存在唯一的 $\mathbb P^1$ 自同构，将 $0,1,\infty$ 映为任意三个互异的点 $p_1,p_2,p_3$。

可以用两种方法计算引力关联数 $\langle H,\tau_1\rangle_{0,1}$。第一种，由伸缩公理，
\[
\langle H,\tau_1\rangle_{0,1}=-\langle H\rangle_{0,1}=-1.
\]
另一种方法是利用除子公理，得到
\[
\langle H,\tau_1\rangle_{0,1}
=1\cdot\langle\tau_1\rangle_{0,1}+\langle H\rangle_{0,1}=-2+1=-1.
\]

接着，利用\hyperref[subsec:gravitational-correlator-properties]{第~\ref*{subsec:gravitational-correlator-properties} 小节}末尾所述的约化过程，计算任意 $d$ 的不变量 $\langle\tau_{2d},1\rangle_{0,d}$ 和 $\langle\tau_{2d-1}H,1\rangle_{0,d}$。为求后者，先注意基本类公理给出
\[
\langle\tau_1H,1\rangle_{0,1}=\langle H\rangle_{0,1}=1.
\]
对一般 $d$，增加三个标记点，并选 $H$ 作为丰沛类。考虑类
\[
c_1(\mathcal L_1)^{2d-1}\cup e_1^*(H)\cup e_3^*(H)\cup e_4^*(H),
\]
将它分别与线性等价关系 $D(2,3\mid4,5)\sim D(2,5\mid3,4)$ 两端在映射 $s:\overline M_{0,5}(\mathbb P^1,d)\to\overline M_{0,4}$ 下的拉回作杯积。然后在 $\overline M_{0,5}(\mathbb P^1,d)$ 上积分并使用分裂公理。可以看到，许多项为零，非零项给出
\begin{equation}\label{eq:projective-line-descendant-splitting}
\begin{aligned}
&\langle\tau_{2d-1}H,1,H,H\rangle_{0,d}\langle1,H,1\rangle_{0,0}\\
&\quad=\langle\tau_{2d-1}H,1,1,1\rangle_{0,d-1}\langle H,H,H\rangle_{0,1}.
\end{aligned}
\end{equation}
由 Gromov--Witten 理论的点映射公理，$\langle1,H,1\rangle_{0,0}=1$；又如上所述，$\langle H,H,H\rangle_{0,1}=1$。于是，除子公理和基本类公理将式~\eqref{eq:projective-line-descendant-splitting} 化为
\[
d^2\langle\tau_{2d-1}H,1\rangle_{0,d}
=\langle\tau_{2d-3}H,1\rangle_{0,d-1}.
\]
结合此前对 $\langle\tau_1H,1\rangle_{0,1}$ 的计算，此递推式给出
\begin{equation}\label{eq:projective-line-descendant-factorial}
\langle\tau_{2d-1}H,1\rangle_{0,d}=\frac1{(d!)^2}.
\end{equation}

转而考虑 $\langle\tau_{2d},1\rangle_{0,d}$，基本类公理给出
\[
\langle\tau_2,1\rangle_{0,1}=\langle\tau_1\rangle_{0,1}=-2.
\]
对一般 $d$，考虑类 $c_1(\mathcal L_1)^{2d}\cup e_3^*(H)\cup e_4^*(H)$，如上将它与线性等价关系作杯积，再积分并使用分裂公理。得到
\[
\begin{aligned}
&\langle\tau_{2d},1,H,H\rangle_{0,d}\langle1,H,1\rangle_{0,0}\\
&\quad=\langle\tau_{2d},1,1,1\rangle_{0,d-1}\langle H,H,H\rangle_{0,1},
\end{aligned}
\]
由此可得
\[
d^2\langle\tau_{2d},1\rangle_{0,d}
+2d\langle\tau_{2d-1}H,1\rangle_{0,d}
=\langle\tau_{2d-2},1\rangle_{0,d-1}.
\]
\SourcePage{331}
这里使用了除子公理和基本类公理。结合 $\langle\tau_2,1\rangle_{0,1}$ 的公式和式~\eqref{eq:projective-line-descendant-factorial}，可解此递推式，得到
\begin{equation}\label{eq:projective-line-descendant-harmonic}
\langle\tau_{2d},1\rangle_{0,d}
=\frac{-2}{(d!)^2}\left(1+\frac12+\cdots+\frac1d\right).
\end{equation}
式~\eqref{eq:projective-line-descendant-factorial} 和~\eqref{eq:projective-line-descendant-harmonic} 的这一计算归功于 Pandharipande。这些公式隐含在 \SourceCite{Givental2} 中，例题~\ref{ex:projective-line-i-j-equality} 将说明这一点。
\end{example}

\SourceExampleNumber{10.1.3}{2}
\begin{example}\label{ex:cy-threefold-gravitational}
设 $V$ 为 Calabi--Yau 三维簇。则虚拟基本类 $[\overline M_{0,n}(V,\beta)]^{\mathrm{virt}}$ 的维数为 $n$。将计算形如 $\langle\tau_d\gamma,1\rangle_{0,\beta}$ 的全部引力关联数。次数公理表明，$\langle\tau_d\gamma,1\rangle_{0,\beta}$ 只有在 $2d+\deg\gamma=4$ 时才可能非零。因此，只需计算 $\langle\tau_2,1\rangle_{0,\beta}$、对除子 $D$ 的 $\langle\tau_1D,1\rangle_{0,\beta}$，以及对曲线类 $C$ 的 $\langle C,1\rangle_{0,\beta}=0$。还可假设 $\beta\ne0$。

回顾\hyperref[subsec:instanton-numbers]{第~\ref*{subsec:instanton-numbers} 小节}，有 Gromov--Witten 不变量
\[
\langle I_{0,0,\beta}\rangle=\sum_{k\mid\beta}n_{\beta/k}k^{-3},
\]
其中 $n_\beta$ 是该小节定义的瞬子数。如\hyperref[chap:gromov-witten]{第七章}中一样，令 $N_\beta=\langle I_{0,0,\beta}\rangle$。

现在计算所需的引力关联数。先考虑 $\langle\tau_2,1\rangle_{0,\beta}$。由基本类公理和伸缩公理，
\[
\langle\tau_2,1\rangle_{0,\beta}
=\langle\tau_1\rangle_{0,\beta}=-2N_\beta.
\]
类似地，若 $D$ 为除子，则
\[
\langle\tau_1D,1\rangle_{0,\beta}
=\langle D\rangle_{0,\beta}
=\left(\int_\beta D\right)N_\beta.
\]
最后，若 $C$ 为曲线类，基本类公理给出 $\langle C,1\rangle_{0,\beta}=0$。将在\hyperref[subsec:cy-threefold-j-function]{第~\ref*{subsec:cy-threefold-j-function} 小节}中使用这些计算。
\end{example}

下面的例子中，虚拟基本类起着关键作用。
\SourceExampleNumber{10.1.3}{3}
\begin{example}\label{ex:genus-one-gravitational-euler}
设 $X$ 为任意复维数 $d$ 的光滑射影簇。将计算 $\langle\tau_1\rangle_{1,0}$。满足 $f_*[C]=0$ 的亏格一稳定映射 $f:(C,p)\to X$ 是常值映射。因此，$(C,p)$ 是带一个标记点的稳定曲线，容易得到同构
\[
\overline M_{1,1}(X,0)\simeq X\times\overline M_{1,1},
\]
将 $f:(C,p)\to X$ 与 $(f(C),(C,p))$ 认同。预期维数为 $1$，而 $\overline M_{1,1}(X,0)$ 的实际维数是 $d+1$，所以超额维数为 $d$。由于 $f^*\Theta_X\simeq\mathcal O_C^d$，有 $h^1(C,f^*\Theta_X)=d$。这说明命题~\ref{prop:excess-obstruction-euler-class} 的假设成立。按该命题的记号，需要计算 $R^1\pi_{2*}e_2^*\Theta_X$ 的第 $d$ 个 Chern 类。由上述同构，容易得到
\[
R^1\pi_{2*}e_2^*\Theta_X
\simeq p_2^*R^1\pi_{2*}\mathcal O_{\overline{\mathcal M}_{1,2}}\otimes p_1^*\Theta_X,
\]
其中 $p_i$ 是 $X\times\overline M_{1,1}$ 的两个投影，而右端的 $\pi_2$ 是通常的忘却第二点的映射 $\overline M_{1,2}\to\overline M_{1,1}$。为简便起见，令 $\mathcal L=R^1\pi_{2*}\mathcal O_{\overline{\mathcal M}_{1,2}}$，这是 $\overline{\mathcal M}_{1,1}$ 上的线丛。于是，命题~\ref{prop:excess-obstruction-euler-class}
\SourcePage{332}
和 Chern 类的标准事实表明，虚拟基本类 $\xi=[\overline M_{1,1}(X,0)]^{\mathrm{virt}}$ 为
\begin{equation}\label{eq:genus-one-constant-virtual-class}
\begin{aligned}
\xi&=c_d(p_2^*\mathcal L\otimes p_1^*\Theta_X)\cap[\overline M_{1,1}(X,0)]\\
&=\bigl(p_1^*c_d(X)+p_2^*c_1(\mathcal L)\cup p_1^*c_{d-1}(X)\bigr)
\cap[\overline M_{1,1}(X,0)],
\end{aligned}
\end{equation}
其中 $c_k(X)=c_k(\Theta_X)$ 是切丛的 Chern 类。

在此情形也能描述线丛 $\mathcal L_1$。若 $\widetilde{\mathcal L}_1$ 为 $\overline{\mathcal M}_{1,1}$ 上相应的丛，则 $\mathcal L_1=p_2^*(\widetilde{\mathcal L}_1)$，因为映射 $\overline M_{1,1}(X,0)\to\overline M_{1,1}$ 不收缩任何曲线。还已知
\[
\int_{\overline M_{1,1}}c_1(\widetilde{\mathcal L}_1)=\frac1{24}
\qquad\text{\SourceCite{Mumford2}}.
\]
现在可以计算
\[
\begin{aligned}
\langle\tau_1\rangle_{1,0}
&=\int_{[\overline M_{1,1}(X,0)]^{\mathrm{virt}}}c_1(\mathcal L_1)\\
&=\int_{X\times\overline M_{1,1}}
p_2^*c_1(\widetilde{\mathcal L}_1)\cup
\bigl(p_1^*c_d(X)+p_2^*c_1(\mathcal L)\cup p_1^*c_{d-1}(X)\bigr)\\
&=\left(\int_{\overline M_{1,1}}c_1(\widetilde{\mathcal L}_1)\right)
\left(\int_Xc_d(X)\right)
=\frac1{24}\int_Xc_d(X).
\end{aligned}
\]
第二个等号使用式~\eqref{eq:genus-one-constant-virtual-class}，第三个等号则来自次数原因所给出的 $c_1(\widetilde{\mathcal L}_1)\cup c_1(\mathcal L)=0$。

文献中通常以另一种方式陈述式~\eqref{eq:genus-one-constant-virtual-class}。由 Serre 对偶，$\mathcal L=R^1\pi_{2*}\mathcal O_{\overline{\mathcal M}_{1,2}}$ 对偶于相对对偶化层 $\omega_2$ 沿 $\pi_2:\overline{\mathcal M}_{1,2}\to\overline{\mathcal M}_{1,1}$ 的推前 $\pi_{2*}\omega_2$。习惯上将该层记为 $\mathcal H$。于是，如果把虚拟基本类看作上同调类，式~\eqref{eq:genus-one-constant-virtual-class} 成为
\[
[\overline M_{1,1}(X,0)]^{\mathrm{virt}}
=p_1^*c_d(X)-p_2^*c_1(\mathcal H)\cup p_1^*c_{d-1}(X).
\]
\SourceCite{BCOV1} 解释了这个公式，\SourceCite{Katz3} 在稳定映射的背景中重新诠释了它；另见 \SourceCite{Getzler1}。

现在重新计算 $\langle\tau_1\rangle_{1,0}$，将得到一个令人意外而又很好的结果。先注意 $\mathcal L_1=p_2^*(\widetilde{\mathcal L}_1)$，故式~\eqref{eq:gravitational-class-definition} 所定义的引力类 $I_{1,1,0}(\tau_1)$ 化为
\[
\begin{aligned}
I_{1,1,0}(\tau_1)
&=PD^{-1}p_{2*}\bigl((c_1(\mathcal L_1)\cup p_1^*(1))\cap\xi\bigr)\\
&=c_1(\widetilde{\mathcal L}_1)\cup PD^{-1}p_{2*}(p_1^*(1)\cap\xi)\\
&=c_1(\widetilde{\mathcal L}_1)\cup I_{1,1,0}(1).
\end{aligned}
\]
第二个等号来自投影公式，而 $I_{1,1,0}(1)$ 是 $H^0(\overline M_{1,1})$ 中的 Gromov--Witten 类。于是，式~\eqref{eq:gravitational-class-integral} 给出
\[
\begin{aligned}
\langle\tau_1\rangle_{1,0}
&=\int_{\overline M_{1,1}}c_1(\widetilde{\mathcal L}_1)\cup I_{1,1,0}(1)\\
&=\deg(I_{1,1,0}(1))\int_{\overline M_{1,1}}c_1(\widetilde{\mathcal L}_1)
=\frac1{24}\deg(I_{1,1,0}(1)),
\end{aligned}
\]
其中 $\deg(I_{1,1,0}(1))$ 使用同构 $\deg:H^0(\overline M_{1,1})\simeq\mathbb C$ 定义，该同构将 $1$ 映为 $1$。

为求 $\deg(I_{1,1,0}(1))$，令 $\psi:\overline M_{0,3}\to\overline M_{1,1}$ 为将前两个标记点粘合以产生亏格一结点曲线的映射。由于 $\overline M_{0,3}$ 是一个点，有
\[
\deg(I_{1,1,0}(1))=\int_{\overline M_{0,3}}\psi^*I_{1,1,0}(1).
\]
\SourcePage{333}
结合上述两个方程，以及\hyperref[chap:gromov-witten]{第七章}的约化公理和点映射公理，得到
\[
\begin{aligned}
\langle\tau_1\rangle_{1,0}
&=\frac1{24}\int_{\overline M_{0,3}}\psi^*I_{1,1,0}(1)\\
&=\frac1{24}\int_{\overline M_{0,3}}\sum_{i,j}g^{ij}I_{0,3}(1,T_i,T_j)
=\frac1{24}\sum_{i,j}g^{ij}g_{ij}.
\end{aligned}
\]
若 $T_i,T_j$ 都是奇类，则 $g^{ij}=-g^{ji}$；若至少有一个是偶类，则 $g^{ij}=g^{ji}$。因此得到 $\langle\tau_1\rangle_{1,0}=\frac1{24}\sum_i(-1)^ib_i(X)$，其中 $b_i(X)$ 是 $X$ 的第 $i$ 个 Betti 数。比较 $\langle\tau_1\rangle_{1,0}$ 的两种计算，就得到熟悉的公式
\[
\int_Xc_d(X)=\sum_{i=0}^{2d}(-1)^ib_i(X).
\]
\end{example}

\subsection{Virasoro 猜想}\label{subsec:virasoro-conjecture}\index{Virasoro猜想@Virasoro 猜想|hyperpage}
本节最后简要介绍一个关于\IndexLocation{indexloc:0459:333}\IndexLocation{indexloc:0473:333}引力关联数的引人关注的猜想。将\emph{引力势}定义为无穷和
\[
\Phi_{\mathrm{grav}}=\sum_{g=0}^{\infty}\kappa^{2g-2}\Phi_g^{\mathrm{grav}},
\]
其中 $\kappa$ 是参数。再定义\emph{配分函数}
\[
Z=e^{\Phi_{\mathrm{grav}}}.
\]
\emph{Virasoro 猜想}断言，存在一个由变量 $t_d^j$ 的形式微分算子组成的 Virasoro 代数（式~\eqref{eq:virasoro-central-extension}），其中这些算子零化 $Z$ \SourceCite{EHX}。

更准确地说，该 Virasoro 代数包含变量 $t_d^j$ 的显式微分算子 $L_n$，$n\ge-1$，满足 $[L_n,L_m]=(n-m)L_{n+m}$。其中，方程 $L_{-1}Z=0$ 来自基本类公理 \SourceCite{Witten3}；由\IndexLocation{indexloc:0461:333}次数公理、\IndexLocation{indexloc:0463:333}除子公理与\IndexLocation{indexloc:0462:333}伸缩公理的组合，可以推出方程 $L_0Z=0$ \SourceCite{Hori}。猜想认为，对所有 $n\ge-1$，均有 $L_nZ=0$。

有时可以将这个微分算子代数扩张为完整 Virasoro 代数 $L_n$、$n\in\mathbb Z$ 的中心扩张。\hyperref[sec:conformal-field-theories]{附录 B.3} 描述了这种中心扩张，尤其见式~\eqref{eq:virasoro-central-extension}。这个表示的中心荷是 $X$ 的 Euler 示性数 $\chi(X)$。猜想并不声称 $n<-1$ 时也有 $L_nZ=0$。事实上，将交换关系 $[L_n,L_{-n}]=2nL_0+(n^3-n)\chi(X)/12$ 与 $Z\ne0$ 结合，在 $\chi(X)\ne0$ 时可知，对任意 $n>1$，$L_nZ=L_{-n}Z=0$ 都不可能同时成立。

当 $X$ 是一个点时，此猜想成为 Witten 猜想 \SourceCite{Witten3}，Kontsevich 已证明它 \SourceCite{Kontsevich1}。Kontsevich 证明中的一个关键思想是，将配分函数表示为\emph{\IndexLocation{indexloc:0760:333}矩阵积分}，即 Hermitian 矩阵空间上的积分；这导出 Virasoro 代数的一种内在作用。在物理学中，有理由认为，更一般的配分函数也可类似地用\emph{\IndexLocation{indexloc:0761:333}矩阵模型}描述，从而导出 Virasoro 代数的自然作用。

此方法已用于预测射影空间的有理曲线和椭圆曲线 Gromov--Witten 不变量。对于 $\mathbb P^2$，它给出有理曲线不变量的式~\eqref{eq:plane-recursion}，以及椭圆曲线不变量的式~\eqref{eq:plane-genus-one-recursion}。
\SourcePage{334}
\SourceCite{EHX} 提出的 Virasoro 代数只有在 $H^*(X)$ 中每个类的 Hodge 型均为 $(p,p)$ 时才能成立（见 \SourceCite{Borisov2}）。在一般情形中，S. Katz 提出了一对经过修改的 Virasoro 代数；\SourceCite{EJX} 表明，它们在算例中给出了正确的预测。陈述这个猜想时，需要将次数公理替换为使用 Hodge 型的类似公理。

\SourceCite{EX} 还表明，Virasoro 猜想加上一个额外假设，可以导出亏格 $g$ 引力后裔之间的关系，称为\emph{拓扑递推关系}（TRR）。这是因为这些关系将含 $\tau_k$ 的引力后裔约化为含 $\tau_{k-1}$ 的后裔，从而得到递推。\SourceCite{EX} 中描述的 TRR 足以确定 $\mathbb P^2$ 的 $g=2$ Gromov--Witten 不变量。E. Getzler 向我们指出，\SourceCite{Getzler1,Getzler2} 的方法可用于给出同一条 TRR，因而证明 \SourceCite{EX} 所得 TRR 的正确性。下面的引理~\ref{lem:genus-zero-topological-recursion} 将给出 $g=0$ 的一条 TRR。

为研究这个猜想，\SourceCite{LiuT} 将 $L_nZ$ 写为 $\bigl(\sum_{g\ge0}\kappa^{2g-2}\Psi_{g,n}\bigr)Z$，其中 $\Psi_{g,n}$ 是变量 $t_d^j$ 和 $q^\beta$ 的形式表达式，而不依赖 $\kappa$。可以证明，$\Psi_{g,n}$ 只依赖亏格 $\le g$ 的引力关联数。\SourceCite{LiuT} 的主要结果是：若 $H^*(X)$ 只有 $(p,p)$ 型，则对所有 $n\ge-1$，有 $\Psi_{0,n}=0$。

\section{Givental 联络}\label{sec:giventals-connection}\index{Givental联络@Givental 联络|hyperpage}
在\hyperref[sec:dubrovin-formalism]{第~\ref*{sec:dubrovin-formalism} 节}中，已经说明，超流形 $H^*(X)$ 上满足命题~\ref{prop:dubrovin-frobenius} 条件的势函数，都自然导出 $H^*(X,\mathbb C)$ 的切丛上的一族形式平坦联络 $\nabla^\lambda$。当势函数为 Gromov--Witten 势时，得到的联络尤其值得研究，它们与\hyperref[sec:a-model-vhs]{第~\ref*{sec:a-model-vhs} 节}所研究的 A 模型联络密切相关。

这里实际研究的是 Dubrovin 联络的一个扭曲版本，称为\emph{Givental 联络} $\nabla^g$。在\hyperref[subsec:giventals-flat-sections]{第~\ref*{subsec:giventals-flat-sections} 小节}中，将说明如何用引力关联数写出一组 $\nabla^g$ 平坦截面的基，并讨论它与\hyperref[chap:quantum-cohomology]{第八章}的 A 模型联络的关系。\hyperref[subsec:equivariant-giventals-connection]{第~\ref*{subsec:equivariant-giventals-connection} 小节}还将解释 Givental 联络中的参数 $\hbar$ 与等变上同调的关系。

本节构造的平坦截面将用于\hyperref[sec:quantum-cohomology-relations]{第~\ref*{sec:quantum-cohomology-relations} 节}中定义 Givental 函数 $J_X$；它又将在该节导出量子微分方程及\IndexLocation{indexloc:1124:334}量子上同调中的关系。

\subsection{平坦截面}\label{subsec:giventals-flat-sections}
给定光滑射影簇 $X$，令 $\Phi$ 表示 Gromov--Witten 势，$*$ 表示大量子乘积。这里不直接使用\hyperref[sec:dubrovin-formalism]{第~\ref*{sec:dubrovin-formalism} 节}所定义的 $\Phi$ 的 Dubrovin 联络，而是遵循 \SourceCite{Givental2}，对它略作修改。如该节一样，将 $\partial/\partial t_i$ 与 $T_i$ 认同，并将联络看作定义在 $H^*(X,\mathbb C)$ 上的平凡上同调丛上。于是，定义\emph{Givental 联络} $\nabla^g$ 为
\begin{equation}\label{eq:giventals-connection}
\nabla^g_{\frac{\partial}{\partial t_i}}\left(\sum_j a_jT_j\right)
=\hbar\sum_j\frac{\partial a_j}{\partial t_i}T_j
-\sum_j a_jT_i*T_j,
\end{equation}
其中 $\hbar$ 为参数。在\hyperref[subsec:equivariant-giventals-connection]{第~\ref*{subsec:equivariant-giventals-connection} 小节}中，将把 $\hbar$ 与 $H^2(B\mathbb C^*)$ 的一个生成元认同，这将在\hyperref[subsec:quintic-mirror-theorem-statement]{第~\ref*{subsec:quintic-mirror-theorem-statement} 小节}中有用。

虽然由于 $\hbar$ 的存在，$\nabla^g$ 并不满足通常联络的 Leibniz 法则，但仍使用联络的标准术语，尤其使用平坦截面的概念。如\hyperref[subsec:a-model-connection]{第~\ref*{subsec:a-model-connection} 小节}所述，Givental 联络与 Dubrovin 联络 $\nabla^\lambda$ 的关系为
\SourcePage{335}
$\nabla^g=\hbar\nabla^{-\hbar^{-1}}$。因此，式~\eqref{eq:giventals-connection} 的平坦截面正是参数 $\lambda=-\hbar^{-1}$ 时 Dubrovin 联络的平坦截面。

对每个指标 $a=0,\ldots,m$，定义形式截面
\begin{equation}\label{eq:giventals-flat-section-series}
s_a=T_a+\sum_{n=0}^{\infty}\hbar^{-(n+1)}\sum_j
\langle\!\langle\tau_nT_a,T_j\rangle\!\rangle_0\,T^j.
\end{equation}
注意，$\langle\!\langle\tau_nT_a,T_j\rangle\!\rangle_0$ 是变量 $t_0,\ldots,t_m$ 的形式幂级数。若以形式符号 $c$ 表示 $c_1(\mathcal L_1)$，则可将式~\eqref{eq:giventals-flat-section-series} 形式地写为
\begin{equation}\label{eq:giventals-flat-section-symbolic}
s_a=T_a+\sum_j\left\langle\!\left\langle\frac{T_a}{\hbar-c},T_j\right\rangle\!\right\rangle_0T^j,
\end{equation}
因为将分母按 $c/\hbar$ 的幂形式展开，再以 $\tau_n$ 替换 $c^n$，就得到式~\eqref{eq:giventals-flat-section-series}。截面 $s_a$ 在这里具有很好的性质。
\begin{proposition}\label{prop:giventals-flat-sections}
截面 $s_a$ 构成 $\nabla^g$ 平坦截面的一组基。由式~\eqref{eq:giventals-connection}，这等价于方程
\[
\hbar\frac{\partial s_a}{\partial t_i}=T_i*s_a,
\qquad a,i=0,\ldots,m,
\]
其中 $*$ 表示大量子乘积。
\end{proposition}

证明之前，需要 \SourceCite{Witten3} 的如下结果。按\hyperref[subsec:virasoro-conjecture]{第~\ref*{subsec:virasoro-conjecture} 小节}的术语，这个引理是一条拓扑递推关系。
\begin{lemma}\label{lem:genus-zero-topological-recursion}\index{拓扑递推关系|hyperpage}
对 $d_1,d_2,d_3\ge0$、$0\le j_1,j_2,j_3\le m$，有
\[
\begin{aligned}
&\langle\!\langle\tau_{d_1+1}T_{j_1},\tau_{d_2}T_{j_2},\tau_{d_3}T_{j_3}\rangle\!\rangle_0\\
&\quad=\sum_a\langle\!\langle\tau_{d_1}T_{j_1},T_a\rangle\!\rangle_0
\langle\!\langle T^a,\tau_{d_2}T_{j_2},\tau_{d_3}T_{j_3}\rangle\!\rangle_0,
\end{aligned}
\]
其中，如往常一样，$T^a$ 是 $T_a$ 相对于杯积的对偶基。
\end{lemma}
\begin{proof}
对 $n\ge3$，考虑从 $\overline M_{0,n}(X,\beta)$ 到一点的映射，记为 $\pi:\overline M_{0,n}(X,\beta)\to\overline M_{0,3}$，以便利用稳定映射空间的函子性质。这个映射将 $(f,C,p_1,\ldots,p_n)$ 送到从 $C$ 删除 $p_4,\ldots,p_n$ 并收缩全部不稳定分支所得的曲线。令 $\mathcal L_1'$ 为 $\overline M_{0,3}$ 上按 $\mathcal L_1$ 的方式定义的（平凡）丛。计算得到
\begin{equation}\label{eq:cotangent-line-three-point-boundary}
c_1(\pi^*\mathcal L_1')
=c_1(\mathcal L_1)-\sum_{K\sqcup L=\{4,\ldots,n\}}
D_{(\{1\}\cup K\mid\{2,3\}\cup L)},
\end{equation}
其中 $D_{(I\mid J)}$ 表示源曲线有两个分支、标记点 $I$ 位于 $f$ 为常值的一个分支、标记点 $J$ 位于另一个分支的稳定映射轨迹的闭包，并对各可能的曲线类分解取和，如图所示：
\begin{figure}[H]
\centering
\includegraphics[width=0.285\textwidth]{gravitational-recursion-boundary.png}
\end{figure}
\SourcePage{336}
现在将式~\eqref{eq:cotangent-line-three-point-boundary} 乘以 $\prod_a\bigl(c_1(\mathcal L_a)^{d_a}\cup e_a^*(T_{j_a})\bigr)$，对虚拟基本类积分，并使用分裂公理，得到
\begin{equation}\label{eq:genus-zero-topological-recursion-coefficients}
\begin{aligned}
&\langle\tau_{d_1+1}T_{j_1},\tau_{d_2}T_{j_2},\tau_{d_3}T_{j_3},
\tau_{d_4}T_{j_4},\ldots,\tau_{d_n}T_{j_n}\rangle_{0,\beta}\\
&=\sum_{\substack{K\sqcup L=\{4,\ldots,n\}\\\beta=\beta_1+\beta_2}}
\sum_a\pm\langle\tau_{d_1}T_{j_1},\tau_{d_{k_1}}T_{j_{k_1}},\ldots,T_a\rangle_{0,\beta_1}\\
&\hspace{5em}\times\langle T^a,\tau_{d_2}T_{j_2},\tau_{d_3}T_{j_3},
\tau_{d_{l_1}}T_{j_{l_1}},\ldots\rangle_{0,\beta_2},
\end{aligned}
\end{equation}
其中 $K=\{k_1,\ldots\}$、$L=\{l_1,\ldots\}$。式~\eqref{eq:genus-zero-topological-recursion-coefficients} 的符号来自方程左右两端奇上同调类次序的可能改变。

按式~\eqref{eq:genus-g-gravitational-coupling} 展开后，容易看出，式~\eqref{eq:genus-zero-topological-recursion-coefficients} 正是所需等式
\[
\begin{aligned}
&\langle\!\langle\tau_{d_1+1}T_{j_1},\tau_{d_2}T_{j_2},\tau_{d_3}T_{j_3}\rangle\!\rangle_0\\
&=\sum_a\langle\!\langle\tau_{d_1}T_{j_1},T_a\rangle\!\rangle_0
\langle\!\langle T^a,\tau_{d_2}T_{j_2},\tau_{d_3}T_{j_3}\rangle\!\rangle_0
\end{aligned}
\]
中的一个系数。注意，导致式~\eqref{eq:genus-zero-topological-recursion-coefficients} 中符号的奇上同调类重新排列，恰好与比较引理右端变量次序和左端变量次序时的重新排列相同。
\end{proof}

现在可证明命题~\ref{prop:giventals-flat-sections}，即证明 $s_a$ 为平坦截面。
\begin{proof}[命题~\ref{prop:giventals-flat-sections} 的证明]
先观察到
\[
\frac{\partial}{\partial t_i}\langle\!\langle\tau_{d_1}\gamma_1,\ldots,\tau_{d_n}\gamma_n\rangle\!\rangle_0
=\langle\!\langle T_i,\tau_{d_1}\gamma_1,\ldots,\tau_{d_n}\gamma_n\rangle\!\rangle_0.
\]
它由式~\eqref{eq:genus-g-gravitational-coupling} 得到，论证类似于引理~\ref{lem:gw-potential-third-derivatives} 的证明。利用此式，从式~\eqref{eq:giventals-flat-section-series} 计算出
\[
\hbar\frac{\partial s_a}{\partial t_i}
=\sum_{n=0}^{\infty}\sum_j\hbar^{-n}
\langle\!\langle T_i,\tau_nT_a,T_j\rangle\!\rangle_0T^j,
\]
而式~\eqref{eq:giventals-flat-section-series} 和~\eqref{eq:big-quantum-product-couplings} 给出
\begin{equation}\label{eq:giventals-section-product}
\begin{aligned}
T_i*s_a
&=\sum_j\langle\!\langle T_i,T_a,T_j\rangle\!\rangle_0T^j\\
&\quad+\sum_{n=0}^{\infty}\sum_{j,c}\hbar^{-(n+1)}
\langle\!\langle\tau_nT_a,T_j\rangle\!\rangle_0
\langle\!\langle T_i,T^j,T_c\rangle\!\rangle_0T^c.
\end{aligned}
\end{equation}
引理~\ref{lem:genus-zero-topological-recursion} 表明，式~\eqref{eq:giventals-section-product} 右端的第二项为
\[
\sum_{n=0}^{\infty}\sum_c\hbar^{-(n+1)}
\langle\!\langle T_i,\tau_{n+1}T_a,T_c\rangle\!\rangle_0T^c.
\]
将它代回式~\eqref{eq:giventals-section-product}，就得到所需结论。
\end{proof}

命题~\ref{prop:giventals-flat-sections} 的一个推论是
\[
\hbar\frac{\partial s_a}{\partial t_0}=s_a,
\]
\SourcePage{337}
因为 $T_0$ 是大量子乘积的单位元。因此，$s_a$ 通过乘法因子 $e^{t_0/\hbar}$ 依赖 $t_0$。鉴于这一点，经常忽略 $t_0$ 的依赖，因为只要插入因子 $e^{t_0/\hbar}$ 就能将它恢复。

以下有时会限制丛及联络 $\nabla^g$ 的底流形。特别地，在\hyperref[sec:quantum-cohomology-relations]{第~\ref*{sec:quantum-cohomology-relations} 节}中，将限制到子空间
\[
M=H^0(X,\mathbb C)\oplus H^2(X,\mathbb C)\subset H^*(X,\mathbb C).
\]
若 $T_1,\ldots,T_r$ 构成 $H^2(X,\mathbb Q)$ 的基，则 $M$ 的变量为 $t_0,t_1,\ldots,t_r$。为看清此时 $\nabla^g$ 的形式，先注意：由\hyperref[subsec:quantum-product-dependence]{第~\ref*{subsec:quantum-product-dependence} 小节}，将大量子乘积限制到 $H^2(X,\mathbb C)$，恰好得到 $*_{\mathrm{small}}$，只需在式~\eqref{eq:genus-g-gravitational-coupling} 中令 $q^\beta=1$，并使用式~\eqref{eq:restricted-big-quantum-product} 附近说明的约定。又因为 $T_0=1$ 同时是大量子上同调与小量子上同调的单位元，故 $\nabla^g$ 在 $M$ 上的限制为
\begin{equation}\label{eq:restricted-giventals-connection}
\nabla^g_{\frac{\partial}{\partial t_i}}\left(\sum_{j=0}^m a_jT_j\right)
=\hbar\sum_{j=0}^m\frac{\partial a_j}{\partial t_i}T_j
-\sum_{j=0}^m a_jT_i*_{\mathrm{small}}T_j,
\end{equation}
其中 $i=0,1,\ldots,r$。特别地，截面 $s_a$ 限制到 $M$ 后满足
\[
\hbar\frac{\partial s_a}{\partial t_i}=T_i*_{\mathrm{small}}s_a,
\qquad a=0,\ldots,m,\quad i=0,\ldots,r.
\]
限制到 $M$ 后，$s_a$ 的公式有很好的简化。
\begin{proposition}\label{prop:restricted-giventals-flat-sections}
令 $\delta=\sum_{i=1}^r t_iT_i$。则 $s_a$ 在 $M$ 上的限制可写为
\[
\begin{aligned}
s_a=e^{t_0/\hbar}\biggl(&e^{\delta/\hbar}\cup T_a\\
&+\sum_{\beta\ne0}\sum_{j=0}^m e^{\int_\beta\delta}
\left\langle\frac{e^{\delta/\hbar}\cup T_a}{\hbar-c},T_j\right\rangle_{0,\beta}T^j\biggr).
\end{aligned}
\]
\end{proposition}
\begin{proof}
前面已注意到，$t_0$ 通过乘法因子 $e^{t_0/\hbar}$ 出现，因此只需证明 $t_0=0$ 的情形。利用 $\delta=\sum_{i=1}^rt_iT_i$，引力耦合~\eqref{eq:genus-g-gravitational-coupling} 成为
\[
\langle\!\langle\tau_nT_a,T_j\rangle\!\rangle_0
=\sum_{k=0}^{\infty}\sum_\beta\frac1{k!}
\langle\tau_nT_a,T_j,(\delta)^k\rangle_{0,\beta},
\]
其中按上述说明，令 $q^\beta=1$。这里使用式~\eqref{eq:genus-g-gravitational-coupling} 后面引入的记号。

由于 $\delta$ 为除子，可以用除子公理简化这个表达式。当 $\beta\ne0$ 时，容易对 $k$ 归纳证明恒等式
\begin{equation}\label{eq:descendant-divisor-binomial}
\langle\tau_nT_a,T_j,(\delta)^k\rangle_{0,\beta}
=\sum_{\mu+\nu=k}\frac{k!}{\mu!\nu!}
\left(\int_\beta\delta\right)^\mu
\langle\tau_{n-\nu}(T_a\cup\delta^\nu),T_j\rangle_{0,\beta},
\end{equation}
其中，当 $n-\nu<0$ 时，将 $\tau_{n-\nu}(T_a\cup\delta^\nu)$ 取为零。对于 $\beta=0$，需要另一个恒等式，因为 $\overline M_{0,2}(X,0)$ 不存在。此时，对 $k\ge1$，利用除子公理得到
\begin{equation}\label{eq:degree-zero-descendant-divisors}
\langle\tau_nT_a,T_j,(\delta)^k\rangle_{0,0}
=\langle\tau_{n-k+1}(T_a\cup\delta^{k-1}),T_j,\delta\rangle_{0,0}.
\end{equation}
若 $k>n+1$，此式显然为零；若 $k=n+1$，则由 Gromov--Witten 不变量的点映射公理，它化为积分 $\int_XT_a\cup T_j\cup\delta^k$。最后，若 $k<n+1$，有同构 $\overline M_{0,3}(X,0)\simeq X$，将 $(f,C,p_1,p_2,p_3)$ 映为 $f(C)\in X$。此外，$\overline M_{0,4}(X,0)\to\overline M_{0,3}(X,0)$ 就是投影 $X\times\overline M_{0,4}\to X$。
\SourcePage{338}
因此容易看出，$\mathcal L_1$ 在 $\overline M_{0,3}(X,0)$ 上平凡，其 Chern 类为零。于是，当 $k<n+1$ 时，式~\eqref{eq:degree-zero-descendant-divisors} 为零，综上可得
\begin{equation}\label{eq:degree-zero-descendant-kronecker}
\langle\tau_nT_a,T_j,(\delta)^k\rangle_{0,0}
=\left(\int_XT_a\cup T_j\cup\delta^k\right)\delta_{k,n+1}.
\end{equation}
利用式~\eqref{eq:descendant-divisor-binomial} 和~\eqref{eq:degree-zero-descendant-kronecker}，$\langle\!\langle\tau_nT_a,T_j\rangle\!\rangle_0$ 的公式简化为
\[
\begin{aligned}
&\frac1{(n+1)!}\int_XT_a\cup T_j\cup\delta^{n+1}\\
&\quad+\sum_{\beta\ne0}\sum_{k=0}^{\infty}\sum_{\mu+\nu=k}
\frac1{\mu!\nu!}\left(\int_\beta\delta\right)^\mu
\langle\tau_{n-\nu}(T_a\cup\delta^\nu),T_j\rangle_{0,\beta}\\
&=\frac1{(n+1)!}\int_XT_a\cup T_j\cup\delta^{n+1}
+\sum_{\beta\ne0}\sum_{\nu=0}^{\infty}
e^{\int_\beta\delta}\frac1{\nu!}
\langle\tau_{n-\nu}(T_a\cup\delta^\nu),T_j\rangle_{0,\beta}.
\end{aligned}
\]
将此表达式代入 $s_a$ 的定义，得到
\begin{equation}\label{eq:restricted-flat-section-expanded}
\begin{aligned}
s_a={}&T_a+\sum_j\left(\sum_{n=0}^{\infty}\hbar^{-(n+1)}
\frac1{(n+1)!}\int_XT_a\cup T_j\cup\delta^{n+1}\right)T^j\\
&+\sum_{\beta\ne0}\sum_j e^{\int_\beta\delta}
\left(\sum_{\nu=0}^{\infty}\sum_{n\ge\nu}\hbar^{-(n+1)}\frac1{\nu!}
\langle\tau_{n-\nu}(T_a\cup\delta^\nu),T_j\rangle_{0,\beta}\right)T^j.
\end{aligned}
\end{equation}
这个式子看起来复杂，但由恒等式
\begin{equation}\label{eq:poincare-dual-basis-reconstruction}
\sum_j\left(\int_X\phi\cup T_j\right)T^j=\phi,
\qquad\phi\in H^*(X,\mathbb C),
\end{equation}
式~\eqref{eq:restricted-flat-section-expanded} 第一行连同 $T_a$ 简化为 $T_a\cup e^{\delta/\hbar}$。第二行括号内的表达式则可令 $k=n-\nu$ 来简化；由于 $\hbar^{-(n+1)}=\hbar^{-(k+1)}\hbar^{-\nu}$，它成为
\[
\sum_{k=0}^{\infty}\hbar^{-(k+1)}
\langle\tau_k(T_a\cup e^{\delta/\hbar}),T_j\rangle_{0,\beta}.
\]
按式~\eqref{eq:giventals-flat-section-symbolic} 引入的\IndexLocation{indexloc:0476:338}符号记法，它等于
\[
\left\langle\frac{T_a\cup e^{\delta/\hbar}}{\hbar-c},T_j\right\rangle_{0,\beta},
\]
由此立即得到命题。
\end{proof}

\subsection{A 模型联络的平坦截面}\label{subsec:a-model-gravitational-flat-sections}
设 $V$ 是复维数 $d$ 的 Calabi--Yau 流形。在\hyperref[subsec:a-model-connection]{第~\ref*{subsec:a-model-connection} 小节}中，按方程
\[
\nabla_{\frac{\partial}{\partial u_j}}(T_k)=T_j*_{\mathrm{small}}T_k,
\qquad1\le j\le r,\quad0\le k\le m,
\]
定义了 A 模型联络，其中 $u_j=\frac1{2\pi i}t_j$。若将 Givental 联络 $\nabla^g$ 限制到 $H^2(V,\mathbb C)$，则式~\eqref{eq:restricted-giventals-connection} 表明，在 $\hbar=-2\pi i$ 时，$-2\pi i\nabla$ 恰好就是 $\nabla^g$。本小节余下部分固定使用这一 $\hbar$ 值。

于是，\hyperref[subsec:giventals-flat-sections]{第~\ref*{subsec:giventals-flat-sections} 小节}中定义的截面 $s_a$ 限制后成为 A 模型联络 $\nabla$ 的平坦截面。为使用这些截面，引入如下通用记号：给定上同调类 $T\in H^*(V,\mathbb C)$，令
\[
s(T)=e^{\delta/\hbar}\cup T
+\sum_{\beta\ne0}\sum_{j=0}^m e^{\int_\beta\delta}
\left\langle\frac{e^{\delta/\hbar}\cup T}{\hbar-c},T_j\right\rangle_{0,\beta}T^j.
\]
\SourcePage{339}
由命题~\ref{prop:restricted-giventals-flat-sections}，$s(T_a)$ 就是截面 $s_a$；容易推出，对任意 $T\in H^*(V,\mathbb C)$，$s(T)$ 都是 $\nabla$ 的形式平坦截面。沿用\hyperref[subsec:quantum-product-dependence]{第~\ref*{subsec:quantum-product-dependence} 小节}的处理，令 $q^\beta=e^{\int_\beta\delta}$，于是可写成
\begin{equation}\label{eq:a-model-flat-section-gravitational}
s(T)=e^{\delta/\hbar}\cup T
+\sum_{\beta\ne0}\sum_{j=0}^m q^\beta
\left\langle\frac{e^{\delta/\hbar}\cup T}{\hbar-c},T_j\right\rangle_{0,\beta}T^j.
\end{equation}
利用这一公式，容易确定 $\nabla$ 的单值性。假设处于\hyperref[subsec:a-model-asymptotics]{第~\ref*{subsec:a-model-asymptotics} 小节}的情形：$T_1,\ldots,T_r$ 是位于 $V$ 的 Kähler 锥闭包中的整数类，变量 $q_j=e^{t_j}=e^{2\pi iu_j}$ 使 $\nabla$ 成为 $(\Delta^*)^r$ 中原点附近的联络，并采用该小节同样的收敛假设。

令 $\mathcal T_j$ 是绕 $(\Delta^*)^r$ 中第 $j$ 个因子一周所给出的\IndexLocation{indexloc:0901:339}\IndexLocation{indexloc:0902:339}单值性变换。
\begin{theorem}\label{thm:a-model-gravitational-monodromy}
单值性 $\mathcal T_j$ 及其对数 $N_j$ 如下作用于 A 模型联络的平坦截面：
\[
\mathcal T_j(s(T))=s(e^{-T_j}\cup T),
\qquad N_j(s(T))=-s(T_j\cup T).
\]
\end{theorem}
\begin{proof}
由于 $q_j=e^{t_j}$，$\mathcal T_j$ 将 $t_j$ 变为 $t_j+2\pi i$。上面已知 $\nabla$ 是 $\hbar=-2\pi i$ 时 Givental 联络的一个倍数，因此可将 $\mathcal T_j$ 的作用写为 $t_j\mapsto t_j-\hbar$。

现在考察此作用对式~\eqref{eq:a-model-flat-section-gravitational} 的影响。因为 $\delta=\sum_{i=1}^rt_iT_i$，$\mathcal T_j$ 将 $\delta$ 变为 $\delta+2\pi iT_j=\delta-\hbar T_j$。它不影响 $q^\beta=e^{\int_\beta\delta}$，因为 $T_j$ 和 $\beta$ 都是整数类。不过，$\delta$ 也出现在 $e^{\delta/\hbar}$ 中，此处 $\mathcal T_j$ 的作用为
\[
e^{\delta/\hbar}\longmapsto
e^{(\delta-\hbar T_j)/\hbar}
=e^{\delta/\hbar-T_j}=e^{\delta/\hbar}\cup e^{-T_j}.
\]
最后一个等号使用了 $T_j$ 的次数为偶数、因此与 $\delta$ 交换这一事实。将它代入式~\eqref{eq:a-model-flat-section-gravitational}，立即得到 $\mathcal T_j(s(T))$ 的公式；取对数即得 $N_j(s(T))$ 的公式。
\end{proof}

定理~\ref{thm:a-model-gravitational-monodromy} 表明，单值性对数 $N_j$ 本质上是与 $-T_j$ 作杯积。更准确地说，$N_j(s(T_a))=-s(T_j\cup T_a)$ 意味着：与 $-T_j$ 作杯积相对于上同调基 $\{T_a\}$ 的矩阵，恰好是 $N_j$ 相对于平坦截面基 $\{s(T_a)\}$ 的矩阵。这是\hyperref[chap:quantum-cohomology]{第八章}中经常使用的事实。

还可用定理~\ref{thm:a-model-gravitational-monodromy} 研究 $\nabla$ 的典范延拓。回顾\hyperref[subsec:limiting-hodge]{第~\ref*{subsec:limiting-hodge} 小节}：由于 $s(T)$ 是 $\nabla$ 的平坦截面，截面
\[
\widetilde s(T)=\exp\left(-\frac1{2\pi i}\sum_j\log(q_j)N_j\right)s(T)
\]
自然延伸为典范延拓上的截面。于是，由定理~\ref{thm:a-model-gravitational-monodromy} 得到以下公式。
\begin{corollary}\label{cor:a-model-canonical-gravitational-section}
对 $T\in H^*(V,\mathbb C)$，截面 $\widetilde s(T)$ 为
\[
\widetilde s(T)=T+\sum_{\beta\ne0}\sum_{j=0}^m q^\beta
\left\langle\frac{T}{\hbar-c},T_j\right\rangle_{0,\beta}T^j.
\]
\end{corollary}
\SourcePage{340}
\begin{proof}
证明很简单。先注意，$\log(q_j)=t_j$、$\hbar=-2\pi i$ 给出
\[
\begin{aligned}
\widetilde s(T)
&=\exp\left(\hbar^{-1}\sum_jt_jN_j\right)s(T)\\
&=s\left(\exp\left(\hbar^{-1}\sum_jt_j(-T_j)\right)\cup T\right)
=s(e^{-\delta/\hbar}\cup T),
\end{aligned}
\]
第二个等号来自 $N_j(s(T))=s((-T_j)\cup T)$。又有 $\delta=\sum_jt_jT_j$。根据式~\eqref{eq:a-model-flat-section-gravitational}，
\[
\begin{aligned}
s(e^{-\delta/\hbar}\cup T)
&=e^{\delta/\hbar}\cup(e^{-\delta/\hbar}\cup T)\\
&\quad+\sum_{\beta\ne0}\sum_{j=0}^m q^\beta
\left\langle\frac{e^{\delta/\hbar}\cup(e^{-\delta/\hbar}\cup T)}{\hbar-c},T_j\right\rangle_{0,\beta}T^j\\
&=T+\sum_{\beta\ne0}\sum_{j=0}^m q^\beta
\left\langle\frac{T}{\hbar-c},T_j\right\rangle_{0,\beta}T^j.
\end{aligned}
\]
结合这些公式，立即得到推论。
\end{proof}

注意，$\widetilde s(T)$ 的公式只包含 $q^\beta=\prod_iq_i^{\int_\beta T_i}$。如往常一样，可假定 $\beta$ 有效；由于 $T_i$ 位于 Kähler 锥闭包中，有 $\int_\beta T_i\ge0$。因此，$\widetilde s(T)$ 全纯延伸到 $0$。

利用上面的推论，现在终于可以证明\hyperref[chap:quantum-cohomology]{第八章}的命题~\ref{prop:a-model-canonical-flat-basis}。将该命题重新陈述如下。
\begin{corollary}\label{cor:a-model-canonical-section-triangular}
对每个齐次类 $T\in H^*(V,\mathbb C)$，截面 $\widetilde s(T)$ 满足
\[
\widetilde s(T)=T+\text{次数更高的项},\qquad\widetilde s(T)(0)=T.
\]
此外，$N_j$ 作用于 $\widetilde s(T_a)$ 的矩阵，等于与 $-T_j$ 作杯积作用于 $T_a$ 的矩阵。
\end{corollary}
\begin{proof}
最后一个断言立即来自定理~\ref{thm:a-model-gravitational-monodromy} 证明后的说明。对于其他断言，先将推论~\ref{cor:a-model-canonical-gravitational-section} 的符号记法展开，得到
\[
\widetilde s(T)=T+\sum_{\beta\ne0}\sum_{j=0}^m\sum_{n=0}^{\infty}
q^\beta\hbar^{-(n+1)}\langle\tau_nT,T_j\rangle_{0,\beta}T^j.
\]
由于 $V$ 为 Calabi--Yau，\IndexLocation{indexloc:0461:340}次数公理意味着，除满足
\[
2n+\deg T+\deg T_j=2d-2,\qquad d=\dim V
\]
的 $j$ 外，所有项均为零。又因 $\deg T^j=2d-\deg T_j$，有 $\deg T^j=\deg T+2n+2$，故上述求和中任何非零项的次数均大于 $\deg T$。

最后，要证明所有 $q_i=0$ 时 $\widetilde s(T)=T$。前面已指出，可以将求和限制到使
\[
q^\beta=\prod_iq_i^{\int_\beta T_i}
\]
各指数均非负的 $\beta\ne0$。因此，只需证明，每个这样的有效 $\beta$ 至少有一个指数为正。否则，就有 $\int_\beta T_i=0$ 对所有 $i$ 成立；因为 $T_i$ 构成 $H^2(V,\mathbb C)$ 的基，这意味着 $\beta$ 与每个 Kähler 类的配对均为零，与非零有效曲线类的正性矛盾。
\end{proof}

\SourcePage{341}
首次在式~\eqref{eq:giventals-flat-section-symbolic} 中引入符号表达式 $T/(\hbar-c)$ 时，它可能显得有些人为。不过，它对定理~\ref{thm:a-model-gravitational-monodromy} 和推论~\ref{cor:a-model-canonical-gravitational-section} 的证明所带来的简化，已充分说明其价值。在\hyperref[sec:quantum-cohomology-relations]{第~\ref*{sec:quantum-cohomology-relations} 节}中，还会经常遇到这个记号。

\hyperref[chap:quantum-cohomology]{第八章}已研究过 Calabi--Yau 三维簇的 A 模型联络的平坦截面。下面比较那里的公式与这里得到的公式。
\SourceExampleNumber{10.2.2}{1}
\begin{example}\label{ex:threefold-flat-section-comparison}
设 $V$ 为 Calabi--Yau 三维簇。在\hyperref[subsec:a-model-asymptotics]{第~\ref*{subsec:a-model-asymptotics} 小节}中，利用式~\eqref{eq:a-model-potential} 的 Gromov--Witten 势
\[
\begin{aligned}
\Phi&=\frac{(2\pi i)^3}{6}\int_V\left(\sum_{j=1}^ru_jT_j\right)^3
+\sum_{\beta\ne0}N_\beta q^\beta\\
&=\frac{(2\pi i)^3}{6}\int_V\left(\sum_{j=1}^ru_jT_j\right)^3+\Phi_{\mathrm{hol}}
\end{aligned}
\]
构造了 $\nabla$ 的平坦截面（如往常一样，$N_\beta=\langle I_{0,0,\beta}\rangle$）。这里限制到偶上同调。在命题~\ref{prop:a-model-canonical-flat-basis} 的证明中，先用式~\eqref{eq:a-model-flat-sections-threefold} 写出平坦截面，再在式~\eqref{eq:a-model-canonical-sections-threefold} 中，对它们施加 $\exp(-\sum_j u_jN_j)=\exp(-\frac1{2\pi i}\sum_j\log(q_j)N_j)$。特别地，得到公式
\begin{equation}\label{eq:threefold-canonical-unit-potential}
\widetilde s_0=T_0-\frac1{(2\pi i)^3}\sum_j
\frac{\partial\Phi_{\mathrm{hol}}}{\partial u_j}T^j
+\frac2{(2\pi i)^3}\Phi_{\mathrm{hol}}T^0.
\end{equation}
其中 $T_0=1$，$T_j$（$1\le j\le r$）构成 $H^2(V,\mathbb C)$ 的基，$T^j\in H^4(V,\mathbb C)$（$1\le j\le r$）及 $T^0\in H^6(V,\mathbb C)$ 为对偶类。

按\hyperref[subsec:a-model-asymptotics]{第~\ref*{subsec:a-model-asymptotics} 小节}的术语，$\widetilde s_0$ 是 $\nabla^c$ 平坦截面，且已证明其在 $0$ 的值为 $T_0$。推论~\ref{cor:a-model-canonical-gravitational-section} 表明，截面 $\widetilde s(T_0)$ 也对 $\nabla^c$ 平坦，而推论~\ref{cor:a-model-canonical-section-triangular} 表明，它在 $0$ 的值也是 $T_0$。因此 $\widetilde s_0=\widetilde s(T_0)$。

为看清它的含义，用上述上同调基写出 $\widetilde s(T_0)$。由于 $T_0=1$，推论~\ref{cor:a-model-canonical-gravitational-section} 的公式成为
\begin{equation}\label{eq:threefold-canonical-unit-descendants}
\begin{aligned}
\widetilde s(T_0)=T_0+\sum_{\beta\ne0}q^\beta\biggl(&
\left\langle\frac1{\hbar-c},T^0\right\rangle_{0,\beta}T_0
+\sum_j\left\langle\frac1{\hbar-c},T^j\right\rangle_{0,\beta}T_j\\
&+\sum_j\left\langle\frac1{\hbar-c},T_j\right\rangle_{0,\beta}T^j
+\left\langle\frac1{\hbar-c},T_0\right\rangle_{0,\beta}T^0\biggr).
\end{aligned}
\end{equation}
比较式~\eqref{eq:threefold-canonical-unit-potential} 和~\eqref{eq:threefold-canonical-unit-descendants} 的系数。对于 $T^0$ 的系数，先注意
\[
\left\langle\frac1{\hbar-c},T_0\right\rangle_{0,\beta}
=\sum_{n=0}^{\infty}\hbar^{-(n+1)}\langle\tau_n,1\rangle_{0,\beta}
=\hbar^{-3}\langle\tau_2,1\rangle_{0,\beta},
\]
因为由次数公理，其余项均为零。比较 $T^0$ 的系数，并使用上述 $\Phi$ 的公式，得到恒等式
\[
\frac2{(2\pi i)^3}N_\beta
=\hbar^{-3}\langle\tau_2,1\rangle_{0,\beta}.
\]
由于 $\hbar=-2\pi i$，有 $\langle\tau_2,1\rangle_{0,\beta}=-2N_\beta$，恰好与例题~\ref{ex:cy-threefold-gravitational} 一致。对于式~\eqref{eq:threefold-canonical-unit-descendants} 求和中的其他项，由次数公理和\IndexLocation{indexloc:0464:341}基本类公理可知，$T_0,T_j$ 的系数均为零。最后，式~\eqref{eq:threefold-canonical-unit-potential} 和~\eqref{eq:threefold-canonical-unit-descendants} 中 $T^j$ 系数相等，等价于
\[
\langle\tau_1,T_j\rangle_{0,\beta}
=-\left(\int_\beta T_j\right)N_\beta,
\]
它来自\IndexLocation{indexloc:0463:341}除子公理与\IndexLocation{indexloc:0462:341}伸缩公理。
\end{example}

\SourcePage{342}
\subsection{参数 \texorpdfstring{$\hbar$}{hbar}}\label{subsec:equivariant-giventals-connection}
本节最后简要解释，参数 $\hbar$ 如何自然地来自等变上同调。以下讨论取自 \SourceCite{Givental2}。

考虑 $\mathbb P^1$ 上的 $\mathbb C^*$ 作用 $t\cdot(x_0,x_1)=(t^{-1}x_0,x_1)$，令 $\hbar\in H^2(B\mathbb C^*)$ 为相应的生成元，故 $H^*(B\mathbb C^*)=\mathbb C[\hbar]$。可通过代换 $\lambda_0\mapsto\hbar$、$\lambda_1\mapsto0$，将它与例题~\ref{ex:projective-equivariant-localization} 中的 $(\mathbb C^*)^2$ 作用比较。于是
\[
H_{\mathbb C^*}^*(\mathbb P^1)\simeq\mathbb C[p,\hbar]/(p(p-\hbar)).
\]
令 $\mathcal R_{\mathbb C^*}=\mathbb C(\hbar)$ 为 $H^*(B\mathbb C^*)$ 的分式域，这是\hyperref[subsec:atiyah-bott-localization]{第~\ref*{subsec:atiyah-bott-localization} 小节}的记号。则 $\{p/\hbar,(\hbar-p)/\hbar\}$ 构成 $H_{\mathbb C^*}^*(\mathbb P^1)\otimes_{H^*(B\mathbb C^*)}\mathcal R_{\mathbb C^*}$ 作为 $\mathcal R_{\mathbb C^*}$ 向量空间的一组基。这组基尤其适合\hyperref[subsec:atiyah-bott-localization]{第~\ref*{subsec:atiyah-bott-localization} 小节}的局部化定理。（注意，也可将 $\mathcal R_{\mathbb C^*}$ 替换为较小的环 $\mathbb C[\hbar,\hbar^{-1}]$。）

\hyperref[chap:mirror-theorem]{第十一章}将使用如下重要思想。给定稳定映射 $f:(\mathbb P^1,p_1,\ldots,p_n)\to X$，令 $\beta=f_*[\mathbb P^1]\in H_2(X,\mathbb Z)$，则 $f$ 的图像自然给出 $\overline M_{0,n}(X\times\mathbb P^1,(\beta,1))$ 的一个元素。此空间相对于 $\overline M_{0,n}(X,\beta)$ 的一个优点是：$\mathbb C^*$ 通过在 $X$ 上的平凡作用和上述 $\mathbb P^1$ 作用，作用于 $X\times\mathbb P^1$。

由于这个 $\mathbb C^*$ 作用，可使用\hyperref[sec:equivariant-gw]{第~\ref*{sec:equivariant-gw} 节}中定义的等变 Gromov--Witten 不变量。为简化叙述，假定 $X$ 光滑且凸。于是 $X\times\mathbb P^1$ 也凸，定义~\ref{def:equivariant-gw} 适用。给定 $\gamma\in H_{\mathbb C^*}^*(X\times\mathbb P^1)$，得到\IndexLocation{indexloc:0536:342}等变不变量
\[
\begin{aligned}
\langle I_{0,n,(\beta,1)}\rangle_{\mathbb C^*}(\gamma,\ldots,\gamma)
&=\int_{\overline M_{0,n}(X\times\mathbb P^1,(\beta,1))_{\mathbb C^*}}
e_1^*(\gamma)\cup\cdots\cup e_n^*(\gamma)\\
&\in H^*(B\mathbb C^*).
\end{aligned}
\]
将它看作分式域 $\mathcal R_{\mathbb C^*}=\mathbb C(\hbar)$ 中的元素较为方便。利用这些不变量，定义等变势函数
\[
\mathcal F^1(\gamma)=\sum_{n=0}^{\infty}\sum_{\beta\in H_2(X)}\frac1{n!}
\langle I_{0,n,(\beta,1)}\rangle_{\mathbb C^*}(\gamma,\ldots,\gamma)q^\beta.
\]
为将它写成形式幂级数，引入一些变量。注意
\[
\begin{aligned}
H_{\mathbb C^*}^*(X\times\mathbb P^1)\otimes_{H^*(B\mathbb C^*)}\mathcal R_{\mathbb C^*}
=H^*(X)\otimes_{\mathbb C}H_{\mathbb C^*}^*(\mathbb P^1)
\otimes_{H^*(B\mathbb C^*)}\mathcal R_{\mathbb C^*}.
\end{aligned}
\]
因此，若 $T_i$ 为 $H^*(X)$ 的一组上同调基，则 $\{T_i\otimes p/\hbar,T_i\otimes(\hbar-p)/\hbar\}$ 是局部化上同调在 $\mathcal R_{\mathbb C^*}$ 上的一组基。然后将
\[
\gamma=\sum_jt_jT_j\otimes p/\hbar
+\sum_i\tau_iT_i\otimes(\hbar-p)/\hbar
\]
代入 $\mathcal F^1$ 的公式。可将 $t_j$、$\tau_i$ 看作取值于 $\mathcal R_{\mathbb C^*}$ 的变量。

用变量 $t_j,\tau_i$ 表示 $\mathcal F^1$ 后，会出现一些很好的性质。先考虑二阶导数
\[
\Phi_{ij}=\frac{\partial^2\mathcal F^1}{\partial\tau_i\partial t_j}.
\]
将 $\Phi=(\Phi_{ij})$ 看作取矩阵值的形式函数。Givental 表明，$\Phi$ 的转置给出 $\nabla^g$ 的\IndexLocation{indexloc:0440:342}平坦截面。
\begin{proposition}\label{prop:graph-space-flat-matrix}
若 $\Phi^*$ 表示 $\Phi$ 的转置，则对所有 $j$，
\[
\hbar\frac{\partial}{\partial t_j}\Phi^*=T_j*\Phi^*,
\]
其中 $s_j$ 是式~(10.15) 所定义的平坦截面。
\end{proposition}
\SourcePage{343}
\begin{proof}
证明类似于命题~\ref{prop:giventals-flat-sections} 的证明，但因为这里处理的是等变 Gromov--Witten 不变量而非引力关联数，稍为简单。详细证明见 \SourceCite{Givental2}。
\end{proof}

$\Phi_{ij}$ 还有一个很好的方面，即可自然地用引力关联数解释。为此，令
\[
\psi_{ij}=g_{ij}+\left\langle\!\left\langle\frac{T_j}{\hbar-c},T_i\right\rangle\!\right\rangle_0.
\]
由式~\eqref{eq:giventals-flat-section-symbolic}，这就是将 $s_j$ 在对偶基中展开后 $T^i$ 的系数，其中，如往常一样，$\{T^i\}$ 为 $\{T_i\}$ 的对偶基。可以将 $\psi_{ij}$ 看作 $t_i$ 和 $\hbar$ 的函数，即 $\psi_{ij}=\psi_{ij}(t_0,\ldots,t_m,\hbar)$。它与上述函数 $\Phi_{ij}$ 的关系如下。
\begin{lemma}\label{lem:graph-space-descendant-factorization}
\[
\Phi_{ij}=-\sum_{k,\ell}\hbar^{-2}
\psi_{ik}(\tau_0,\ldots,\tau_m,-\hbar)g^{k\ell}
\psi_{j\ell}(t_0,\ldots,t_m,\hbar).
\]
\end{lemma}
\begin{proof}
\SourceCite{Givental2} 证明了这个引理，这里只简述证明。关键是，如\hyperref[chap:localization]{第九章}中一样，用局部化计算左端。$X\times\mathbb P^1$ 上的 $\mathbb C^*$ 作用诱导稳定映射空间上的作用。可以看到，固定点轨迹由稳定映射 $f:C\to X\times\mathbb P^1$ 组成，其中 $C=C'\cup C_0\cup C_\infty$，满足 $f(C_0)\subset X\times\{0\}$、$f(C_\infty)\subset X\times\{\infty\}$，而对某点 $x\in X$，有 $f(C')\subset\{x\}\times\mathbb P^1$。分支 $C_0$、$C_\infty$ 可为空。分支 $C'$ 分别在一个点与 $C_0$、$C_\infty$ 相交，将引理~\ref{lem:node-smoothing-tangent-tensor} 应用于这两个结点，便得到 $\psi_{ij}$ 定义中的 $\hbar-c$ 和 $-\hbar-c$ 因子。
\end{proof}

最后应说明，这些思想也可用于给出命题~\ref{prop:giventals-flat-sections} 的另一证明：再次使用局部化，从引理~\ref{lem:graph-space-descendant-factorization} 和命题~\ref{prop:graph-space-flat-matrix} 推出该结果。

\section{量子上同调中的关系}\label{sec:quantum-cohomology-relations}
本节用式~\eqref{eq:giventals-flat-section-series} 中的平坦截面 $s_a$ 研究量子上同调。将 $\nabla^g$ 看作底空间 $M=H^0(X,\mathbb C)\oplus H^2(X,\mathbb C)$ 上平凡丛 $M\times H^*(X,\mathbb C)$ 的形式联络。因此，变量为 $t_0,t_1,\ldots,t_r$，$\nabla^g$ 由式~\eqref{eq:restricted-giventals-connection} 中的小量子乘积给出。

\subsection{Givental 函数 \texorpdfstring{$J$}{J}}\label{subsec:giventals-j-function}\index{Givental $J$ 函数|hyperpage}\index{J函数@$J$ 函数|hyperpage}
遵循 Givental，定义函数
\begin{equation}\label{eq:giventals-j-pairing-definition}
J=\sum_j\langle s_j,1\rangle T^j,
\end{equation}
其中 $\langle\alpha,\beta\rangle=\int_X\alpha\cup\beta$ 是上同调上通常的交配对。利用式~\eqref{eq:giventals-flat-section-series}，也可将 $J$ 的定义写成
\begin{equation}\label{eq:giventals-j-descendant-series}
J=J(t_0,\ldots,t_m,\hbar^{-1})
=1+\sum_{n=0}^{\infty}\sum_{a=0}^m\hbar^{-(n+1)}
\langle\!\langle\tau_nT_a,1\rangle\!\rangle_0T^a.
\end{equation}
\hyperref[subsec:quantum-product-dependence]{第~\ref*{subsec:quantum-product-dependence} 小节}已说明，变量 $t_1,\ldots,t_r$ 在\IndexLocation{indexloc:1122:343}大量子上同调中以指数 $q_j=e^{t_j}$ 出现。因此，将小量子乘积看作大量子乘积的限制时，给出 $*_{\mathrm{small}}$ 的级数是
\SourcePage{344}
$q_j$ 的形式幂级数。为简便起见，以下将 $*_{\mathrm{small}}$ 简记为 $*$。

设 $P(\hbar\partial/\partial t,e^t,\hbar)$ 是以下量的形式幂级数：
\[
\hbar\frac{\partial}{\partial t_0},\ldots,
\hbar\frac{\partial}{\partial t_r},e^{t_0},\ldots,e^{t_r},\hbar.
\]
为明确起见，可以且将假定，指数项都放在微分算子左侧，尽管下面的定理~\ref{thm:quantum-differential-relations} 并不需要这一条件。

若 $H^*(X,\widehat C)$ 表示系数为 $\widehat C=\mathbb C[[q_1,\ldots,q_r]]$ 的上同调，则以 $P(T,q,0)$ 表示对 $P(\hbar\partial/\partial t,e^t,\hbar)$ 作代换
\[
\hbar\frac{\partial}{\partial t_j}\longmapsto T_j,
\qquad e^{t_j}\longmapsto q_j,\qquad\hbar\longmapsto0
\]
并将微分算子的复合替换为量子乘积后，所得的 $H^*(X,\widehat C)$ 中的形式幂级数。若 $P$ 还含 $e^{t_0}$，则在系数环中相应增添标量参数 $q_0$。关键结果如下。
\begin{theorem}\label{thm:quantum-differential-relations}
设 $P(\hbar\partial/\partial t,e^t,\hbar)J=0$，其中 $P$ 是按上述记号表示的形式幂级数。则关系 $P(T,q,0)=0$ 在小量子上同调中成立。
\end{theorem}
\begin{proof}
以下证明由 B. Kim 提出。引入\emph{对偶 Givental 联络} $\widetilde\nabla^g$，定义为
\[
\widetilde\nabla^g_{\frac{\partial}{\partial t_i}}\left(\sum_j a_jT_j\right)
=\hbar\sum_j\frac{\partial a_j}{\partial t_i}T_j
+\sum_j a_jT_i*T_j.
\]
简单计算可验证，对变量 $t_0,\ldots,t_m$ 的取值于 $H^*(X,\mathbb C)$ 的函数 $G,H$，有恒等式
\begin{equation}\label{eq:dual-giventals-pairing-derivative}
\hbar\frac{\partial}{\partial t_i}\langle G,H\rangle
=\left\langle\nabla^g_{\frac{\partial}{\partial t_i}}G,H\right\rangle
+\left\langle G,\widetilde\nabla^g_{\frac{\partial}{\partial t_i}}H\right\rangle.
\end{equation}
又可算得
\begin{equation}\label{eq:dual-giventals-classical-symbol}
\widetilde\nabla^g_{\frac{\partial}{\partial t_{i_1}}}\cdots
\widetilde\nabla^g_{\frac{\partial}{\partial t_{i_k}}}1
=T_{i_1}*\cdots*T_{i_k}+\hbar A(T,q,\hbar),
\end{equation}
其中 $A$ 是变量 $T_i,q_i,\hbar$ 的形式幂级数。

表达式 $PJ$ 定义为 $\sum_jP\langle s_j,1\rangle T^j$。由于 $T^0,\ldots,T^m$ 线性无关，$PJ=0$ 意味着，对所有 $j$，均有 $P(\hbar\partial/\partial t,e^t,\hbar)\langle s_j,1\rangle=0$。命题~\ref{prop:giventals-flat-sections} 又给出 $\nabla^g_{\partial/\partial t_i}s_j=0$。重复应用式~\eqref{eq:dual-giventals-pairing-derivative} 和~\eqref{eq:dual-giventals-classical-symbol}，可得
\[
\begin{aligned}
0=P\langle s_j,1\rangle
&=\langle s_j,P(T,q,0)\rangle
+\hbar\langle s_j,B(T,q,\hbar)\rangle,\\[-2pt]
&\hspace{10em}j=0,\ldots,m,
\end{aligned}
\]
其中 $B(T,q,\hbar)$ 是形式幂级数。由于 $s_j$ 构成一组基，且交配对非退化，先得到 $P(T,q,0)+\hbar B(T,q,\hbar)=0$。这是 $\hbar$ 的非负幂级数恒等式，因此可令 $\hbar=0$，得到 $P(T,q,0)=0$，并有
\[
\langle s_j,P(T,q,0)\rangle=0,\qquad j=0,\ldots,m.
\]
所需关系便在 $H^*(X,\widehat C)$ 的小量子乘积下成立。
\end{proof}

这个定理导出如下定义。
\begin{definition}\label{def:quantum-differential-operator}\index{量子微分方程|hyperpage}
满足定理~\ref{thm:quantum-differential-relations} 假设的微分算子 $P(\hbar\partial/\partial t,e^t,\hbar)$ 称为\emph{\IndexLocation{indexloc:1151:344}\IndexLocation{indexloc:1154:344}量子微分算子}。方程 $P(\hbar\partial/\partial t,e^t,\hbar)Y=0$（由 $Y=J$ 满足）称为\emph{量子微分方程}。
\end{definition}

\SourcePage{345}
也可以如下理解定理~\ref{thm:quantum-differential-relations}。令 $P_{\widetilde\nabla^g}$ 表示对 $P(\hbar\partial/\partial t,e^t,\hbar)$ 作代换
\[
\hbar\frac{\partial}{\partial t_i}\longmapsto
\widetilde\nabla^g_{\frac{\partial}{\partial t_i}}
\]
所得的表达式。由于 $s_j$ 是 $\nabla^g$ 的平坦截面，式~\eqref{eq:dual-giventals-pairing-derivative} 给出
\begin{equation}\label{eq:j-annihilator-dual-connection}
P\langle s_j,1\rangle=\langle s_j,P_{\widetilde\nabla^g}1\rangle.
\end{equation}
因此，对于 Givental 函数 $J$，有等价关系
\begin{equation}\label{eq:j-annihilator-equivalences}
PJ=0\ \Longleftrightarrow\ P\langle s_j,1\rangle=0\text{ 对所有 }j
\ \Longleftrightarrow\ P_{\widetilde\nabla^g}1=0.
\end{equation}
第一个等价关系来自定理~\ref{thm:quantum-differential-relations} 的证明；第二个来自式~\eqref{eq:j-annihilator-dual-connection}，因为 $s_j$ 成基。这解释了式~\eqref{eq:giventals-j-pairing-definition} 中“$1$”的作用。

注意，若 $P(T,q)=0$ 是小量子上同调中的关系，则 $P(\hbar\partial/\partial t,e^t)$ 不一定是量子微分算子（见例题~\ref{ex:cy-quantum-equation-relation}）。不过，若此关系的次数至多为二，且指数项放在微分算子的左边，就容易看出，$P(\hbar\partial/\partial t,e^t)$ 是量子微分算子。此外，下面的例题~\ref{ex:second-order-quantum-relations} 将说明，某些特殊情形下，可直接由量子上同调中的关系导出量子微分方程。最后还应说明，\hyperref[chap:mirror-theorem]{第十一章}的例题~\ref{ex:toric-fano-and-hirzebruch-quantum-ring} 将利用定理~\ref{thm:quantum-differential-relations} 计算某些环面簇的量子上同调。

接着要简化 $J$ 的公式，以便计算量子微分算子的一些例子。以下引理记录了两个经常使用的公式。
\begin{lemma}\label{lem:giventals-j-two-formulas}
$X$ 的 Givental $J$ 函数由以下两个公式给出：
\[
\begin{aligned}
J&=e^{(t_0+\delta)/\hbar}\left(1+\sum_{\beta\ne0}\sum_{a=0}^m q^\beta
\left\langle\frac{T_a}{\hbar-c},1\right\rangle_{0,\beta}T^a\right)\\
&=e^{(t_0+\delta)/\hbar}\left(1+\sum_{\beta\ne0}q^\beta PD^{-1}e_{1*}
\left(\frac1{\hbar-c}\cap[\overline M_{0,2}(X,\beta)]^{\mathrm{virt}}\right)\right),
\end{aligned}
\]
其中 $\delta=\sum_{i=1}^rt_iT_i$，$q^\beta=e^{\int_\beta\delta}$。第二个等式中，$PD$ 是 Poincaré 对偶，$e_1:\overline M_{0,2}(X,\beta)\to X$ 是第一个标记点处的赋值映射。
\end{lemma}
\begin{proof}
回顾 $J$ 是用平坦截面 $s_a$ 构造的，而我们限制在 $M=H^0(X,\mathbb C)\oplus H^2(X,\mathbb C)$ 上，因此命题~\ref{prop:restricted-giventals-flat-sections} 给出
\[
s_a=e^{t_0/\hbar}\left(e^{\delta/\hbar}\cup T_a
+\sum_{\beta\ne0}\sum_{j=0}^m q^\beta
\left\langle\frac{e^{\delta/\hbar}\cup T_a}{\hbar-c},T_j\right\rangle_{0,\beta}T^j\right).
\]
\SourcePage{346}
故 Givental 函数 $J$ 为
\begin{equation}\label{eq:j-function-flat-section-expansion}
\begin{aligned}
J&=\sum_{a=0}^m\langle s_a,1\rangle T^a\\
&=e^{t_0/\hbar}\biggl(\sum_{a=0}^m\langle e^{\delta/\hbar}\cup T_a,1\rangle T^a\\
&\hspace{4em}+\sum_{a=0}^m\sum_{\beta\ne0}\sum_{j=0}^m q^\beta
\left\langle\frac{e^{\delta/\hbar}\cup T_a}{\hbar-c},T_j\right\rangle_{0,\beta}
\langle T^j,1\rangle T^a\biggr).
\end{aligned}
\end{equation}
第二行的求和容易简化：
\[
\sum_{a=0}^m\langle e^{\delta/\hbar}\cup T_a,1\rangle T^a
=\sum_{a=0}^m\langle e^{\delta/\hbar},T_a\rangle T^a=e^{\delta/\hbar},
\]
最后一个等号来自式~\eqref{eq:poincare-dual-basis-reconstruction}。第三行更简单，因为
\[
\langle T^j,1\rangle=
\begin{cases}1,&T_j=T_0=1,\\0,&\text{其他情形}.
\end{cases}
\]
故第三行只有 $j=0$ 的项有贡献。因此
\begin{equation}\label{eq:j-function-exponential-insertion}
J=e^{t_0/\hbar}\left(e^{\delta/\hbar}
+\sum_{\beta\ne0}\sum_{a=0}^m q^\beta
\left\langle\frac{e^{\delta/\hbar}\cup T_a}{\hbar-c},1\right\rangle_{0,\beta}T^a\right).
\end{equation}
还可进一步简化。虽然上同调基 $\{T_a\}$ 及其对偶基 $\{T^a\}$ 出现在式~\eqref{eq:j-function-exponential-insertion} 中，$J$ 并不依赖所用的基，故可用任意其他基替换 $\{T_a\}$。甚至可取非齐次基；由于使用形式系数，可取基 $\{e^{-\delta/\hbar}\cup T_a\}$。它的对偶基为 $\{e^{\delta/\hbar}\cup T^a\}$，因此
\[
\begin{aligned}
J=e^{t_0/\hbar}\biggl(&e^{\delta/\hbar}\\
&+\sum_{\beta\ne0}\sum_{a=0}^m q^\beta
\left\langle\frac{e^{\delta/\hbar}\cup(e^{-\delta/\hbar}\cup T_a)}{\hbar-c},1\right\rangle_{0,\beta}
e^{\delta/\hbar}\cup T^a\biggr),
\end{aligned}
\]
简化后便是所需的第一个公式。

为证明第二个公式，注意
\begin{equation}\label{eq:descendant-j-evaluation-integral}
\left\langle\frac{T_a}{\hbar-c},1\right\rangle_{0,\beta}
=\int_{[\overline M_{0,2}(X,\beta)]^{\mathrm{virt}}}\frac{e_1^*(T_a)}{\hbar-c}.
\end{equation}
由 $e_1$ 的投影公式，此式等于
\[
\int_XT_a\cup PD^{-1}\left(e_{1*}
\left(\frac1{\hbar-c}\cap[\overline M_{0,2}(X,\beta)]^{\mathrm{virt}}\right)\right).
\]
将 $\phi=PD^{-1}\bigl(e_{1*}((1/(\hbar-c))\cap[\overline M_{0,2}(X,\beta)]^{\mathrm{virt}})\bigr)$ 代入式~\eqref{eq:poincare-dual-basis-reconstruction}，再用式~\eqref{eq:descendant-j-evaluation-integral}，得到
\[
\begin{aligned}
\sum_{a=0}^m\left\langle\frac{T_a}{\hbar-c},1\right\rangle_{0,\beta}T^a
=PD^{-1}\left(e_{1*}\left(\frac1{\hbar-c}\cap[\overline M_{0,2}(X,\beta)]^{\mathrm{virt}}\right)\right).
\end{aligned}
\]
这表明，引理中的第二个公式立即来自第一个。
\end{proof}

\SourcePage{347}
引理~\ref{lem:giventals-j-two-formulas} 的一个重要推论，结合基本类公理所给出的 $\langle T_a,1\rangle_{0,\beta}=0$，是
\begin{equation}\label{eq:j-function-asymptotic-normalization}
J=e^{(t_0+\delta)/\hbar}\bigl(1+o(\hbar^{-1})\bigr),
\end{equation}
如 \SourceCite{Givental4} 所指出的。这个事实对\hyperref[sec:givental-mirror-approach]{第~\ref*{sec:givental-mirror-approach} 节}中定理~\ref{thm:givental-projective-mirror} 的陈述至关重要。

下面给出两个例题，既给出 $J$ 的公式，也说明量子微分方程如何导出量子上同调中的关系。第~\ref{subsec:cy-threefold-j-function} 小节还将给出更多例题。
\SourceExampleNumber{10.3.1}{1}
\begin{example}\label{ex:projective-line-j-function}
计算 $X=\mathbb P^1$ 的 $J$。若 $H$ 是超平面截面类（这里就是点类），则 $\delta=t_1H$；对于 $\beta=d\in H_2(\mathbb P^1,\mathbb Z)=\mathbb Z$，有 $\int_\beta\delta=dt_1$。由引理~\ref{lem:giventals-j-two-formulas}，
\[
J=e^{(t_0+t_1H)/\hbar}\left(1+\sum_{d=1}^{\infty}q^d
\left(\left\langle\frac H{\hbar-c},1\right\rangle_{0,d}1
+\left\langle\frac1{\hbar-c},1\right\rangle_{0,d}H\right)\right).
\]
利用式~\eqref{eq:projective-line-descendant-factorial} 和~\eqref{eq:projective-line-descendant-harmonic}，它成为
\[
\begin{aligned}
J=e^{(t_0+t_1H)/\hbar}\biggl(1+\sum_{d=1}^{\infty}q^d\biggl(&
\frac{\hbar^{-2d}}{(d!)^2}1\\
&-\frac{2\hbar^{-(2d+1)}}{(d!)^2}
\left(1+\frac12+\cdots+\frac1d\right)H\biggr)\biggr)\\
=e^{(t_0+t_1H)/\hbar}\biggl(1+\sum_{d=1}^{\infty}
&\left(\frac q{\hbar^2}\right)^d\frac1{(d!)^2}
\left(1-2\frac H\hbar\left(1+\frac12+\cdots+\frac1d\right)\right)\biggr).
\end{aligned}
\]
直接验证可知，$(\hbar\,d/dt_1)^2-e^{t_1}$ 零化 $J$，因而是量子微分算子。于是，定理~\ref{thm:quantum-differential-relations} 在代换 $\hbar\,d/dt_1\mapsto H$、$e^{t_1}\mapsto q$ 后，给出小量子上同调关系 $H^2=q$，与\hyperref[subsec:a-model-correlation]{第~\ref*{subsec:a-model-correlation} 小节}的例题~\ref{ex:projective-quantum-ring} 一致。

还可如下得到一个很有趣的 $J$ 公式。由于 $H^2=0$，上式可写成
\[
J=e^{(t_0+t_1H)/\hbar}\sum_{d=0}^{\infty}q^d
\frac1{\left(d!\hbar^{d-1}\left(\sum_{j=1}^d1/j\right)H+d!\hbar^d\right)^2},
\]
再利用一次 $H^2=0$，得到
\[
J=e^{(t_0+t_1H)/\hbar}\sum_{d=0}^{\infty}q^d
\frac1{\bigl((H+\hbar)(H+2\hbar)\cdots(H+d\hbar)\bigr)^2}.
\]
\hyperref[subsec:givental-nef-projective-mirror]{第~\ref*{subsec:givental-nef-projective-mirror} 小节}将说明，这是 Givental 版本的镜像定理的一个特例。
\end{example}

\SourceExampleNumber{10.3.1}{2}
\begin{example}\label{ex:second-order-quantum-relations}
考虑 $n$ 维射影空间 $X=\mathbb P^n$，令 $H$ 为 $\mathbb P^n$ 的超平面类。取上同调基 $1,H,\ldots,H^n$，其中变量 $t_0,t_1$ 对应于类 $1,H$。令 $q=e^{t_1}$，则由例题~\ref{ex:projective-quantum-ring}，小量子乘积满足 $H^{n+1}=e^{t_1}$。

证明这个关系导出一个量子微分方程。由于当 $i\le n$ 时，量子幂 $H^i$ 不依赖任何 $t_0,t_1$，推导式~\eqref{eq:dual-giventals-classical-symbol} 的分析在此给出更精确的恒等式
\[
\left(\widetilde\nabla^g_{\frac d{dt_1}}\right)^{n+1}1=H^{n+1}.
\]
\SourcePage{348}
又因小量子上同调中 $H^{n+1}=e^{t_1}$，此式成为
\[
\left(\widetilde\nabla^g_{\frac d{dt_1}}\right)^{n+1}1=e^{t_1}.
\]
由式~\eqref{eq:j-annihilator-equivalences}，得到
\begin{equation}\label{eq:projective-space-quantum-operator}
\left(\hbar\frac d{dt_1}\right)^{n+1}-e^{t_1}
\end{equation}
是量子微分算子。

也可反过来推导。命题~\ref{prop:projective-space-i-equals-j} 断言，$\mathbb P^n$ 的 Givental 函数 $J$ 由下式给出：
\begin{equation}\label{eq:projective-space-j-hypergeometric}\index{Givental $J$ 函数!射影空间|hyperpage}
J=e^{(t_0+t_1H)/\hbar}\sum_{d=0}^{\infty}e^{dt_1}
\frac1{\bigl((H+\hbar)(H+2\hbar)\cdots(H+d\hbar)\bigr)^{n+1}}.
\end{equation}
目前先注意这个公式的两点。第一，当 $n=1$ 时，它与例题~\ref{ex:projective-line-j-function} 的计算一致（令 $q=e^{t_1}$ 即可）；第二，对于一般的 $n$，可证明式~\eqref{eq:projective-space-j-hypergeometric} 右端被式~\eqref{eq:projective-space-quantum-operator} 的微分算子消去。因此，\hyperref[chap:mirror-theorem]{第十一章}证明式~\eqref{eq:projective-space-j-hypergeometric} 后，即可推出式~\eqref{eq:projective-space-quantum-operator} 是量子微分算子；再由定理~\ref{thm:quantum-differential-relations}，得到 $\mathbb P^n$ 小量子上同调中的关系 $H^{n+1}=e^{t_1}=q$。
\end{example}

\subsection{Calabi--Yau 三维簇}\label{subsec:cy-threefold-j-function}
本节余下部分均假设 $V$ 是 Calabi--Yau 三维簇。先解释如何用\hyperref[subsec:instanton-numbers]{第~\ref*{subsec:instanton-numbers} 小节}所讨论的 Gromov--Witten 不变量
\[
N_\beta=\langle I_{0,0,\beta}\rangle
\]
表示 Givental $J$ 函数。如例题~\ref{ex:threefold-flat-section-comparison} 一样，限制到偶上同调，将 $H^{\mathrm{even}}(V)$ 的上同调基排列成 $T_0=1$，$T_1,\ldots,T_r$ 生成 $H^2(V)$，而 $T^a\in H^4(V)$（$1\le a\le r$）与 $T^0\in H^6(V)$ 为对偶类。还令 $\delta=\sum_{i=1}^rt_iT_i$。

现在给出 $J$ 的公式。引理~\ref{lem:giventals-j-two-formulas} 给出
\[
J=e^{(t_0+\delta)/\hbar}\left(1+\sum_{\beta\ne0}q^\beta
\left(\hbar^{-2}\sum_{a=1}^r\langle\tau_1T_a,1\rangle_{0,\beta}T^a
+\hbar^{-3}\langle\tau_2,1\rangle_{0,\beta}T^0\right)\right).
\]
由次数公理与基本类公理容易得到此式。在例题~\ref{ex:cy-threefold-gravitational} 中，已证明
\[
\langle\tau_1T_a,1\rangle_{0,\beta}=\left(\int_\beta T_a\right)N_\beta,
\qquad\langle\tau_2,1\rangle_{0,\beta}=-2N_\beta,
\]
故上述公式化简为
\begin{equation}\label{eq:cy-threefold-j-invariants}
J=e^{(t_0+\delta)/\hbar}\left(1+\sum_{\beta\ne0}q^\beta
\left(\hbar^{-2}N_\beta\sum_{a=1}^r\left(\int_\beta T_a\right)T^a
-2\hbar^{-3}N_\beta T^0\right)\right).
\end{equation}
不过，式~\eqref{eq:poincare-dual-basis-reconstruction} 容易推出，$\beta\in H_2(V,\mathbb C)$ 的 Poincaré 对偶为 $\sum_{a=1}^r(\int_\beta T_a)T^a\in H^4(V,\mathbb C)$。略微滥用记号，仍用 $\beta$ 表示这个 $H^4(V,\mathbb C)$ 中的类，就可将 Calabi--Yau 三维簇的 Givental 函数公式改写为
\begin{equation}\label{eq:cy-threefold-j-poincare-dual}
J=e^{(t_0+\delta)/\hbar}\left(1+\hbar^{-2}\sum_{\beta\ne0}N_\beta q^\beta\beta
-2\hbar^{-3}\sum_{\beta\ne0}N_\beta q^\beta\,pt\right).
\end{equation}
\SourcePage{349}
若有独立的方法计算 $J$，则与上述公式比较系数，便可计算 $N_\beta$。Givental 在 \SourceCite{Givental4} 中断言，环面完全交具有这样的计算方法。\hyperref[subsec:givental-nef-projective-mirror]{第~\ref*{subsec:givental-nef-projective-mirror} 小节}的例题~\ref{ex:givental-quintic-mirror-potential} 将以五次三维簇说明这种方法。

下面解释 Givental $J$ 函数与 Gromov--Witten 势及小量子乘积的关系。结果如下。
\begin{proposition}\label{prop:cy-threefold-j-potential}\index{Givental $J$ 函数!Calabi--Yau 三维簇|hyperpage}
若 $\Phi$ 是 Calabi--Yau 三维簇 $V$ 的 Gromov--Witten 势，则 $V$ 的 Givental $J$ 函数为
\[
\begin{aligned}
J=e^{t_0/\hbar}\biggl(1+\hbar^{-1}\sum_{a=1}^rt_aT_a
&+\hbar^{-2}\sum_{a=1}^r\frac{\partial\Phi}{\partial t_a}T^a\\
&+\hbar^{-3}\left(\sum_{a=1}^rt_a\frac{\partial\Phi}{\partial t_a}-2\Phi\right)T^0\biggr).
\end{aligned}
\]
此外，对任意 $1\le i,j\le r$，有
\[
\frac{\partial^2J}{\partial t_i\partial t_j}
=\hbar^{-2}e^{(t_0+\delta)/\hbar}T_i*T_j.
\]
\end{proposition}
\begin{proof}
回顾式~\eqref{eq:cy-total-gw-potential}，Gromov--Witten 势为
\[
\Phi=\frac16\int_V\delta^3+\sum_{\beta\ne0}N_\beta q^\beta
=\frac16\int_V\delta^3+\Phi_{\mathrm{hol}},\qquad\delta=\sum_{a=1}^rt_aT_a.
\]
利用此式及式~\eqref{eq:cy-threefold-j-invariants}，容易得到
\[
J=e^{(t_0+\delta)/\hbar}\left(1+\hbar^{-2}\sum_{a=1}^r
\frac{\partial\Phi_{\mathrm{hol}}}{\partial t_a}T^a-2\hbar^{-3}\Phi_{\mathrm{hol}}T^0\right).
\]
将乘积展开，得到
\[
\begin{aligned}
J=e^{t_0/\hbar}\biggl(&1+\hbar^{-1}\delta
+\frac{\hbar^{-2}}2\delta^2+\frac{\hbar^{-3}}6\delta^3\\
&+\hbar^{-2}\sum_{a=1}^r\frac{\partial\Phi_{\mathrm{hol}}}{\partial t_a}T^a
+\hbar^{-3}\delta\sum_{a=1}^r\frac{\partial\Phi_{\mathrm{hol}}}{\partial t_a}T^a
-2\hbar^{-3}\Phi_{\mathrm{hol}}T^0\biggr).
\end{aligned}
\]
由此出发，利用以下恒等式即可得到所需的 $J$ 公式，细节留给读者：
\begin{equation}\label{eq:cy-threefold-potential-euler-identities}
\begin{aligned}
\delta\sum_{a=1}^r\frac{\partial\Phi_{\mathrm{hol}}}{\partial t_a}T^a
&=\left(\sum_{a=1}^rt_a\frac{\partial\Phi_{\mathrm{hol}}}{\partial t_a}\right)T^0,\\
\frac16\delta^3&=\frac16\left(\int_V\delta^3\right)T^0
=\frac13\sum_{a=1}^rt_a\frac\partial{\partial t_a}
\left(\frac16\int_V\delta^3\right)T^0,\\
\frac12\delta^2&=\frac12\sum_{a=1}^r\left(\int_V\delta^2\cup T_a\right)T^a
=\sum_{a=1}^r\frac\partial{\partial t_a}
\left(\frac16\int_V\delta^3\right)T^a.
\end{aligned}
\end{equation}
注意，第二行使用 Euler 公式，第三行使用式~\eqref{eq:poincare-dual-basis-reconstruction}。

为证明 $\partial^2J/\partial t_i\partial t_j$ 的公式，将命题陈述中的 $J$ 公式对 $t_i,t_j$ 求导，容易得到
\[
\begin{aligned}
\frac{\partial^2J}{\partial t_i\partial t_j}=e^{t_0/\hbar}\biggl(&
\hbar^{-2}\sum_{a=1}^r\frac{\partial^3\Phi}{\partial t_i\partial t_j\partial t_a}T^a\\
&+\hbar^{-3}\left(\sum_{a=1}^rt_a\frac{\partial^3\Phi}{\partial t_i\partial t_j\partial t_a}\right)T^0\biggr).
\end{aligned}
\]
\SourcePage{350}
利用类似于式~\eqref{eq:cy-threefold-potential-euler-identities} 第一行的恒等式，得到
\[
\begin{aligned}
\frac{\partial^2J}{\partial t_i\partial t_j}
&=\hbar^{-2}e^{(t_0+\delta)/\hbar}\sum_a
\frac{\partial^3\Phi}{\partial t_i\partial t_j\partial t_a}T^a\\
&=\hbar^{-2}e^{(t_0+\delta)/\hbar}T_i*T_j.
\end{aligned}
\]
第二个等号来自命题~\ref{prop:giventals-flat-sections} 证明后的讨论，那里解释了：小量子乘积就是大量子乘积在 $M=H^0(V)\oplus H^2(V)$ 上的限制。
\end{proof}

应用此命题，计算一个有趣的量子微分算子。
\SourceExampleNumber{10.3.2}{1}
\begin{example}\label{ex:cy-quantum-equation-relation}
设 Calabi--Yau 三维簇 $V$ 满足 $h^2(V)=1$。令 $H$ 为 $H^2(V)$ 的生成元，$C\in H^4(V)$ 为对偶生成元，用 $pt\in H^6(V,\mathbb Z)$ 表示点类，故 $H\cup C=pt$。如往常一样，变量 $t_0,t_1$ 对应于 $1,H$，并令 $q=e^{t_1}$。

\hyperref[subsec:a-model-hodge-filtration]{第~\ref*{subsec:a-model-hodge-filtration} 小节}研究了三点函数 $Y(q)=\langle H,H,H\rangle$，并在式~\eqref{eq:one-parameter-quantum-yukawa} 中指出关系 $H*H=Y(q)C$。利用命题~\ref{prop:cy-threefold-j-potential}，得到
\[
\left(\hbar\frac d{dt_1}\right)^2J
=e^{(t_0+t_1H)/\hbar}H*H=e^{(t_0+t_1H)/\hbar}Y(q)C.
\]
容易由此推出
\[
\begin{aligned}
\left(\hbar\frac d{dt_1}\right)^2
\left(\frac{(\hbar\,d/dt_1)^2J}{Y(q)}\right)
&=\left(\hbar\frac d{dt_1}\right)^2e^{(t_0+t_1H)/\hbar}C\\
&=e^{(t_0+t_1H)/\hbar}H\cup H\cup C=0,
\end{aligned}
\]
因为 $H^8(V)=0$。因此
\begin{equation}\label{eq:cy-threefold-quantum-operator}
\left(\hbar\frac d{dt_1}\right)^2\circ\frac1{Y(q)}\circ
\left(\hbar\frac d{dt_1}\right)^2
\end{equation}
是量子微分算子。由定理~\ref{thm:quantum-differential-relations}，它在 $V$ 的\IndexLocation{indexloc:1139:350}小量子上同调环中给出关系 $H*H*(H*H/Y)=0$。

由于 $H^8(V)=0$，这个关系本身相当平凡。不过，微分方程~\eqref{eq:cy-threefold-quantum-operator} 仍很有趣，主要是因为它与\hyperref[sec:yukawa-couplings]{第~\ref*{sec:yukawa-couplings} 节}为 Calabi--Yau 三维簇的 Yukawa 耦合所推导的 Picard--Fuchs 方程 \eqref{eq:normalized-yukawa-pf} 极其相似。这并非巧合，而是镜像对称的一个关键部分。下面的定理~\ref{thm:picard-fuchs-quantum-equivalence} 将进一步研究 Picard--Fuchs 方程与量子微分方程的关系。

还可用这个例题说明\hyperref[subsec:giventals-j-function]{第~\ref*{subsec:giventals-j-function} 小节}的评论：定理~\ref{thm:quantum-differential-relations} 的逆命题可能不成立，即量子上同调中的关系并非都给出量子微分方程。对于这里的 Calabi--Yau 三维簇 $V$，量子上同调中确有关系 $H*H*H*H=0$，
\SourcePage{351}
但 $(\hbar\,d/dt_1)^4J\ne0$（当 $Y$ 非常数时）。这是因为，用式~\eqref{eq:dual-giventals-pairing-derivative} 作简单计算，得到
\[
\left(\hbar\frac d{dt_1}\right)^4\langle s_j,1\rangle
=\hbar^2\frac{d^2Y}{dt_1^2}\langle s_j,C\rangle
+2\hbar\frac{dY}{dt_1}\langle s_j,pt\rangle.
\]
由于 $C$ 和 $pt$ 是线性无关的上同调类，且 $Y$ 非常数，容易推出 $(\hbar\,d/dt_1)^4\langle s_j,1\rangle$ 不可能对所有 $j$ 都为零。
\end{example}

本节最后解释，定理~\ref{thm:quantum-differential-relations} 与例题~\ref{ex:threefold-flat-section-comparison} 中讨论的 A 模型联络 $\nabla$ 有何关系。回顾，自然变量为 $u_j$，它与 $t_j$ 的关系是 $t_j=2\pi iu_j$。考虑算子
\[
P_\nabla=P\left(\nabla_{\frac\partial{\partial u_1}},\ldots,
\nabla_{\frac\partial{\partial u_r}},q_1,\ldots,q_r\right),
\qquad q_j=e^{2\pi iu_j}.
\]
假定 $P_\nabla$ 是联络项的多项式，其系数为 $q_j$ 的形式幂级数，并假定所有 $q_j$ 都放在联络项的左边。

\begingroup\linespread{1.18}\selectfont
算子 $P_\nabla$ 可作用于定义 $\nabla$ 的丛 $\mathcal H$ 的任意截面。特别地，有由 $1=T_0$ 给出的截面。若 $P_\nabla1=0$，则称 $P_\nabla$ 为 \emph{Picard--Fuchs 算子}。

这个术语源于镜像对称，因为在\hyperref[sec:mirror-symmetry-revisited]{第~\ref*{sec:mirror-symmetry-revisited} 节}讨论的镜像猜想中，截面 $1$ 对应于镜像上的归一化三形式 $\widetilde\Omega$，而通常的 Picard--Fuchs 算子恰好是消去 $\widetilde\Omega$ 的算子。

给定 $P_\nabla$，将它写成
\begin{equation}\label{eq:a-model-picard-fuchs-operator}
P_\nabla=\sum_\alpha A_\alpha(q)\nabla^\alpha,
\end{equation}
其中 $\alpha=(a_1,\ldots,a_r)$ 为多重指标，
\[
\nabla^\alpha=\nabla_{\frac\partial{\partial u_1}}^{a_1}\cdots
\nabla_{\frac\partial{\partial u_r}}^{a_r}.
\]
如往常一样，令 $|\alpha|=\sum_i a_i$，并以 $m=\max_\alpha\{|\alpha|:A_\alpha(q)\ne0\}$ 表示 $P_\nabla$ 的阶数。定义 $P_\nabla$ 的\emph{$\hbar$ 齐次化}为微分算子
\begin{equation}\label{eq:picard-fuchs-homogenization}
P(\hbar\partial/\partial t,e^t,\hbar)
=\sum_\alpha\left(\frac\hbar{2\pi i}\right)^{m-|\alpha|}
A_\alpha(e^t)\left(\hbar\frac\partial{\partial t}\right)^\alpha,
\end{equation}
其中 $(\hbar\partial/\partial t)^\alpha$ 的含义显而易见。（下面会看到为何需要 $2\pi i$。）还令
\[
P_m(T,q)=\sum_{|\alpha|=m}A_\alpha(q)T^\alpha,
\]
其中 $T^\alpha$ 是用 $T_1,\ldots,T_r$ 在 $H^*(V)$ 的小量子乘积下组成的乘积。可将 $P_m(T,q)$ 看作算子 $P_\nabla$ 的主部（特征形式）的量子版本。

这些对象的关系如下。
\begin{theorem}\label{thm:picard-fuchs-quantum-equivalence}\index{量子微分方程!Picard--Fuchs 算子|hyperpage}
设 $V$ 为 Calabi--Yau 三维簇，$P_\nabla$ 为式~\eqref{eq:a-model-picard-fuchs-operator} 中的 $m$ 阶算子。若 $P(\hbar\partial/\partial t,e^t,\hbar)$ 由式~\eqref{eq:picard-fuchs-homogenization} 定义，则：
\SourcePage{352}
\begin{enumerate}[label=(\roman*)]
\item $P_\nabla$ 是 Picard--Fuchs 算子，当且仅当 $P(\hbar\partial/\partial t,e^t,\hbar)$ 是量子微分算子。换言之，
\[
P_\nabla1=0\quad\Longleftrightarrow\quad
P(\hbar\partial/\partial t,e^t,\hbar)J=0.
\]
\item 若 $P_\nabla$ 是 Picard--Fuchs 算子，则在小量子上同调中 $P_m(T,q)=0$。
\end{enumerate}
\end{theorem}
\begin{proof}
先指出，此时量子微分方程具有一个特殊性质。令 $J$ 如式~\eqref{eq:cy-threefold-j-poincare-dual}，将 $\hbar$ 视为变量；以 $J_0$ 表示把其中的 $\hbar$ 替换为非零复数 $\hbar_0$ 后所得的表达式。断言
\begin{equation}\label{eq:quantum-operator-hbar-specialization}
P(\hbar\partial/\partial t,e^t,\hbar)J=0
\quad\Longleftrightarrow\quad
P(\hbar_0\partial/\partial t,e^t,\hbar_0)J_0=0.
\end{equation}
\looseness=-1 为说明这一点，先注意 $P(\hbar\partial/\partial t,e^t,\hbar)$ 关于 $\hbar$ 齐次，因此它等于 $\hbar^m\widetilde P$，其中算子 $\widetilde P$ 不含 $\hbar$。此外，$\hbar$ 在 $J$ 的系数中也是齐次出现的。将式~\eqref{eq:cy-threefold-j-invariants} 展开，并暂时除去不受这些 $t_1,\ldots,t_r$ 微分算子影响的公因子 $e^{t_0/\hbar}$，可写成
\[
J=1+\hbar^{-1}\sum_{j=1}^rA_jT_j
+\hbar^{-2}\sum_{j=1}^rB_jT^j+\hbar^{-3}CT^0,
\]
其中 $A_j,B_j,C$ 为 $q_1,\ldots,q_r$ 的函数。由 $P=\hbar^m\widetilde P$，$PJ=0$ 等价于
\[
\widetilde P1=\widetilde PA_j=\widetilde PB_j=\widetilde PC=0
\qquad\text{对所有 }j.
\]
由于 $\hbar_0\ne0$，这也等价于 $P(\hbar_0\partial/\partial t,e^t,\hbar_0)J_0=0$，从而得到式~\eqref{eq:quantum-operator-hbar-specialization}。

\looseness=-1 现在证明定理。由式~\eqref{eq:quantum-operator-hbar-specialization}，可将 $\hbar$ 特化为任意非零常数。故在证明余下部分中，不失一般性，取 $\hbar=2\pi i$。马上就会看到这是正确的选择。

由式~\eqref{eq:j-annihilator-equivalences}，已知
\begin{equation}\label{eq:quantum-operator-dual-annihilation}
P(\hbar\partial/\partial t,e^t,\hbar)J=0
\quad\Longleftrightarrow\quad P_{\widetilde\nabla^g}1=0,
\end{equation}
其中 $\widetilde\nabla^g$ 为对偶 Givental 联络，$P_{\widetilde\nabla^g}$ 由代换 $\hbar\partial/\partial t_j\mapsto\widetilde\nabla^g_{\partial/\partial t_j}$ 定义。

关键观察是：由于 $\hbar=2\pi i$，对偶 Givental 联络 $\widetilde\nabla^g$ 恰好是 $2\pi i\nabla$。因此
\[
\widetilde\nabla^g_{\frac\partial{\partial t_j}}
=2\pi i\nabla_{\frac\partial{\partial t_j}}
=\nabla_{\frac\partial{\partial u_j}},
\]
第二个等号使用了 $t_j=2\pi iu_j$。由于 $\hbar=2\pi i$，比较式~\eqref{eq:a-model-picard-fuchs-operator} 和~\eqref{eq:picard-fuchs-homogenization}，可见 $P_{\widetilde\nabla^g}$ 恰好为 $P_\nabla$。于是式~\eqref{eq:quantum-operator-dual-annihilation} 给出
\[
P(\hbar\partial/\partial t,e^t,\hbar)J=0\quad\Longleftrightarrow\quad P_\nabla1=0.
\]
这证明了第一部分，第二部分立即由定理~\ref{thm:quantum-differential-relations} 得到。
\end{proof}

\enlargethispage{2\baselineskip}
\looseness=-1 这个定理表明，Calabi--Yau 三维簇的 A 模型联络与小量子上同调有着深刻联系。不过，镜像对称所预言的、它与镜像的 B 模型之间的联系更深。下一章将证明镜像定理的几个特殊情形，最终将一直以来所研究的这些美妙数学联系起来。
\par\endgroup
